% =====================================================================
%  Capitolo 11 — Il vocabolario del level designer: blockout e i suoi
%  cinque strumenti, coordinate, gerarchie, snapping, pivot, mesh e
%  collisioni, luci, e il primo livello costruito da zero in Unreal
%  Engine (primo anno)
%  Fonte: trascrizione della lezione GD11, che è la seconda lezione del
%  corso di Level Design del primo anno (docente: Stefano Zocchi, lo
%  stesso del capitolo 8), di cui la lezione è la diretta prosecuzione.
%  Contenuti, nell'ordine della lezione: la struttura delle lezioni del
%  primo anno; la revisione dei blockout della propria stanza (Francesca,
%  Antonino, Federica, Mattia, Francesco, Alberto, Valerio), con le
%  digressioni su Win+Shift+S, il "pizzico di caos" e il verosimile; lo
%  Starter Content assente in Unreal 5.8; lavoro artistico e lavoro di
%  gameplay; il blockout (Call of Duty 2019, Uncharted, il gameplay che
%  regna sovrano); i cinque strumenti del blockout (primitivi, kit
%  modulari, brush e modellazione, sculpting del terreno, spline); che
%  cosa si farà nel primo anno; asset e actor; coordinate globali e
%  locali (alla lavagna); la pausa; lo spazio delle coordinate
%  nell'editor; genitori, figli e gruppi; Ctrl+Z e la selezione; lo
%  snapping (griglia, rotazione, campi numerici); gli assi di rotazione;
%  il Content Browser e la lingua dell'editor; il pivot; mesh e
%  collisioni; le luci (point, spot, rect, directional); l'importazione
%  manuale dello Starter Content; il livello da zero (piano, luce
%  direzionale, Lumen, luce ambientale e cielo, la cartella Lighting);
%  materiali e texture; le modalità di visualizzazione; il post process
%  volume; il task (una stanza creativa); le domande e le presenze.
%  Integrazioni dal libro "An Architectural Approach to Level Design"
%  (Christopher W. Totten), segnalate come approfondimenti.
%
%  Pacchetti richiesti nel preambolo del libro (già presenti nel
%  main.tex): \usepackage{booktabs}, \usepackage{amsmath},
%  \usepackage{tikz} con
%    \usetikzlibrary{shapes.geometric,shapes.symbols,shapes.misc,
%      arrows.meta,positioning,calc,fit}
%  (le librerie vanno caricate nel preambolo, non nel capitolo).
%  Tutte le etichette di questo capitolo hanno il prefisso "gd11:"
%  (gd11:sec:..., gd11:fig:..., gd11:tab:...) e tutti gli stili TikZ il
%  prefisso "gd11", per non collidere con gli altri capitoli.
%
%  Nota sulla fonte: la trascrizione automatica è buona nella parte
%  teorica, ma molto rovinata altrove. L'inizio della registrazione, le
%  due pause e i momenti di lavoro individuale sono stati trascritti
%  come frasi senza senso o ripetute in loop ("Il numero di persone è
%  molto più basso dopo il raggio di persone", "ti sanno, ti sanno",
%  "vado a rio", "e questo no"), con brani captati da una fonte esterna
%  (frasi sul "mercato interno dell'Unione Europea", sulla "patente",
%  sugli "alzheimeri"): sono artefatti, omessi. Durante la pausa la
%  registrazione cattura anche conversazioni private (sulla musica, su
%  Jimi Hendrix e sulla chitarra suonata al contrario) e, a fine
%  lezione, un breve scambio privato del docente con uno studente di un
%  altro anno su una consegna: non fanno parte della lezione e non sono
%  riportati. Molti nomi propri (di giochi, persone, strumenti) sono
%  storpiati ("Arial", "il Real", "Ruben", "Elder Ring", "AlphaFlash"):
%  dove il riferimento è ricostruibile con certezza lo si riporta nella
%  forma corretta; dove non lo è, lo si segnala in nota invece di
%  indovinare. Le figure e le tabelle sono ricostruzioni schematiche di
%  ciò che il docente mostra e descrive (anche alla lavagna), non le
%  slide originali.
% =====================================================================

\chapter[Blockout, coordinate e luci]{Il vocabolario del level designer:
blockout, coordinate, pivot, mesh e luci}
\label{gd11:cap:blockout-luci}

Il capitolo~\ref{gd8:cap:leveldesign} ha aperto il corso di \textbf{Level
Design} del primo anno: chi è il level designer, perché ``level design'' è un
termine ombrello, e i primissimi passi nell'editor di Unreal Engine, con il
task di ricostruire \textbf{la propria stanza} usando solo forme di base. Questa
lezione, tenuta dallo stesso docente (Stefano Zocchi), ne è la diretta
prosecuzione: è la ``\textbf{seconda parte del ciclo di lezioni
introduttive}''. Le lezioni introduttive, spiega il docente, servono a vedere
``un po' di concetti, un po' di altra roba di base'', con momenti di
spiegazione e momenti di pratica.

Il docente lo dice subito con franchezza: questa è ``una delle lezioni più
frontali'' dell'anno, ``una lunga lista di termini che è necessario che
sappiate se vogliamo iniziare a parlare con cognizione di causa degli
argomenti del nostro corso''. Sono concetti legati all'uso di Unreal, ma anche
\textbf{concetti di base del mondo dello sviluppo di videogiochi e del level
design}: molti studenti, almeno in parte, li conoscono già, ma il docente
sceglie di vederli ``in una sezione formale all'interno del corso'' per
assicurarsi che tutti sappiano esattamente che cosa sono. ``La lezione di oggi
sarà un po' noiosa, o magari queste cose vi appassionano, non lo so.''

La lezione è anche segnata da un imprevisto tecnico: il computer dell'aula è
\textbf{nuovo}, il docente non aveva ancora usato \textbf{Unreal Engine 5.8}
(per lavoro usa Unity o altre versioni di Unreal), e nella 5.8 Epic ha cambiato
diverse cose, tra cui il fatto che lo \textbf{Starter Content} non è più
incluso. Una parte della lezione è quindi spesa a risolvere il problema, e il
capitolo riporta anche la procedura trovata, perché è quella che gli studenti
devono seguire.

La struttura della lezione, e quindi del capitolo, è questa:
\begin{enumerate}
  \item la \textbf{struttura delle lezioni} del primo anno e la
        \textbf{revisione dei blockout} della propria stanza, con le
        digressioni sugli screenshot, sul ``pizzico di caos'' e sul
        \textbf{verosimile} (sezioni~\ref{gd11:sec:struttura}
        e~\ref{gd11:sec:revisione});
  \item l'introduzione alla lezione e lo Starter Content mancante
        (sezione~\ref{gd11:sec:intro});
  \item la distinzione fondamentale tra \textbf{lavoro artistico} e
        \textbf{lavoro di gameplay} nel level design
        (sezione~\ref{gd11:sec:arte-gameplay});
  \item il \textbf{blockout}: che cos'è, perché si fa, gli esempi di
        \emph{Call of Duty} e \emph{Uncharted}, e il gameplay che ``regna
        sovrano'' (sezioni~\ref{gd11:sec:blockout}
        e~\ref{gd11:sec:esempi-blockout});
  \item i \textbf{cinque strumenti} del blockout: primitivi, kit modulari,
        brush e modellazione, scultura del terreno, spline
        (sezioni~\ref{gd11:sec:strumenti}--\ref{gd11:sec:spline}), e che
        cosa si farà nel primo anno (sezione~\ref{gd11:sec:primo-anno});
  \item la terminologia dell'editor: \textbf{asset e actor},
        \textbf{coordinate globali e locali}, \textbf{genitori e figli},
        \textbf{gruppi}, Ctrl+Z, \textbf{snapping}, assi di rotazione,
        Content Browser, \textbf{pivot}
        (sezioni~\ref{gd11:sec:asset}--\ref{gd11:sec:pivot});
  \item \textbf{mesh e collisioni} (sezione~\ref{gd11:sec:mesh});
  \item le \textbf{luci}: point light, spot light, rect light e
        directional light (sezione~\ref{gd11:sec:luci});
  \item l'importazione manuale dello \textbf{Starter Content}
        (sezione~\ref{gd11:sec:starter});
  \item il \textbf{livello costruito da zero}: piano, sole, luce ambientale,
        cielo, la cartella \emph{Lighting}
        (sezione~\ref{gd11:sec:da-zero});
  \item \textbf{materiali e texture}, le \textbf{modalità di
        visualizzazione} e il \textbf{post process volume}
        (sezioni~\ref{gd11:sec:materiali}--\ref{gd11:sec:postprocess});
  \item il \textbf{task} per la lezione successiva, le domande e la chiusura
        (sezioni~\ref{gd11:sec:task} e~\ref{gd11:sec:chiusura}).
\end{enumerate}

% ---------------------------------------------------------------------
\section{La struttura delle lezioni del primo anno}
\label{gd11:sec:struttura}

\subsection{Cinque minuti per chi arriva}

Come fa sempre, il docente lascia qualche minuto ``per dare il tempo a un
pochino di gente di arrivare''. La lezione è ibrida: la maggior parte della
classe è in aula, mentre \textbf{Mattia} segue da casa (e combatterà a lungo con
l'audio di Teams). Il docente sa che una persona gli ha scritto per dire che
sarebbe stata assente; l'appello vero e proprio verrà fatto ``da lezione'',
cioè alla fine, per essere sicuri che ci siano tutti quelli che ci devono
essere (sezione~\ref{gd11:sec:chiusura}).

\subsection{Prima i task, tutti insieme}

Il docente spiega come sono organizzate le sue lezioni, ``soprattutto il primo
anno'': \textbf{la prima parte della lezione è dedicata a guardare in classe,
tutti insieme, i task} che gli studenti hanno fatto.
\begin{itemize}
  \item Questa abitudine \textbf{svanisce progressivamente} nel corso
        dell'anno, e dipende anche da \textbf{quante persone ci sono in
        classe}: con quattro o cinque persone c'è il tempo di vedere tutto
        insieme, con dieci o quindici molto meno.
  \item Il primo anno, però, il docente ci tiene a cominciare le lezioni
        \textbf{correggendo} o, come in questo caso, \textbf{dando
        un'occhiata}: le prime lezioni sono di pratica, e non è troppo
        rigido nel correggere.
  \item Lo scopo è doppio: tutti possono guardare i lavori degli altri e
        \textbf{rendersi conto del livello generale}; e, dato che ci saranno
        molti casi in cui tutti fanno \textbf{lo stesso errore}, il docente
        non deve ripeterlo più volte: ``posso parlarne di uno e lo vedete
        tutti''.
\end{itemize}

% ---------------------------------------------------------------------
\section{La revisione dei blockout della stanza}
\label{gd11:sec:revisione}

``Per oggi c'era un task. Il task era?'' La classe risponde subito (``ottimo,
dieci su dieci''): il \textbf{blockout della propria stanza}, il primo task di
level design (sezione~\ref{gd8:sec:task}). Si procede \textbf{in ordine
alfabetico}, proiettando gli screenshot caricati nella cartella dei task.

\subsection{Francesca: cubi, piani e un CD sotto la sedia}

La prima è Francesca (il docente si fa correggere la pronuncia del cognome).
``Bella stanza.'' Ci sono delle \textbf{decorazioni}: la studentessa spiega che
sono stilizzate, e il docente approva (``ci può stare''). Fa notare un
elemento storto, ma il giudizio complessivo è ``bel lavoro''. Alla domanda se
abbia usato solo cubi, Francesca risponde di aver usato \textbf{cubi e piani}
(i piani per i poster), e c'è persino un dettaglio sotto la sedia (``dovrebbe
essere un CD'').

\paragraph{Lo scopo dell'esercizio.} ``Per questo task di solito non ho niente
da dire'': il docente lo guarda soprattutto per assicurarsi che tutti
l'abbiano fatto. Il lavoro è ``molto preciso'', ``decisamente molto
dettagliato''; ma lo scopo dell'esercizio, ricorda, è \textbf{prendere la mano
con gli strumenti}: passare dallo spostamento alla rotazione alla scala e così
via, in modo che diventi \textbf{istintivo}, ``non ci dovete più pensare una
volta che l'avete imparato''. Per questo va benissimo; e non si perde troppo
tempo su ogni lavoro.

\subsection{Antonino: una PlayStation 5, e la cartella condivisa}

La stanza di Antonino è ``molto più ristretta''. Il docente si diverte a
riconoscere gli oggetti: ``fammi indovinare, questa è una PlayStation 5''; e
poi il case del computer.

\paragraph{Un account condiviso.} Nel frattempo, il docente spiega perché
sullo schermo compare il nome di un altro docente (Mattia Podini): sta usando
l'\textbf{account generale} che è stato lasciato sul computer dell'aula.

\paragraph{Non si cancella nella cartella dei task.} Nella cartella compare
anche una ``bozza''. Regola pratica: gli studenti \textbf{non possono
cancellare} i file caricati nella cartella condivisa. Se serve eliminare
qualcosa, bisogna \textbf{scrivere ``da cancellare''} (nel nome del file): il
docente lo cancellerà appena lo vede.

\subsection{Federica: abbiamo appena traslocato}

Federica mostra ``solo la versione migliore'', e ha già usato anche i
\textbf{materiali}: ``siamo già più avanti''. La stanza è ``un po' spoglia,
forse''; ma c'è una ragione narrativa: ``abbiamo appena traslocato, quindi la
stanza è vuota''. ``Questo è un nuovo modo per imbrogliare un task'', scherza
il docente; ma il lavoro è buono, e soprattutto contiene un dettaglio che vale
per tutti.

\paragraph{Un pizzico di caos.} Le \textbf{scatole} del trasloco sono messe
\textbf{sghembe}, di traverso, e questo crea una visuale più interessante.
Stessa cosa per il \textbf{libro sul comodino}: ``se l'avessi messo
rigidamente regolare, sarebbe stato molto strano, tra virgolette, da vedere''.
\begin{quote}
Mettere un pizzico di caos in quello che facciamo rende gli ambienti più
interessanti, rende gli ambienti più \textbf{vissuti}.
\end{quote}
``Anche queste sono cose a cui arriveremo'': è un primo assaggio del lavoro di
\emph{set dressing} (sezione~\ref{gd8:sec:cosa-fa}).

\subsection{Mattia: le foto allo schermo, e uno spadone}

Mattia Maestri è il primo, nell'esperienza del docente, a consegnare
\textbf{fotografie dello schermo} invece che screenshot. Segue da casa, ma
all'inizio non lo si sente: il microfono è disattivato, poi l'audio della
televisione dell'aula non funziona (``questo computer è nuovo, tra l'altro, e
questo potrebbe crearci un paio di problemi dopo''), finché Mattia non si
collega anche dal telefono, scoprendo nel frattempo di avere le cuffie rotte.
La spiegazione delle foto: ha dovuto formattare da poco il PC fisso e non ha
ancora reinstallato gli strumenti per fare e mandare gli screenshot.

\paragraph{Un trucco su Windows.} Il docente ne approfitta per insegnare a tutti
una cosa ``abbastanza utile'':
\begin{itemize}
  \item \textbf{Win+Shift+S} è la scorciatoia di Windows per fare uno
        \textbf{screenshot}: si seleziona una parte dello schermo, e
        l'immagine viene salvata \textbf{negli appunti}; con \textbf{Ctrl+V}
        la si incolla dove serve. Funziona ovunque in Windows.
  \item Esiste anche una funzione \textbf{interna a Unreal} che permette di
        prendere uno screenshot direttamente dall'immagine renderizzata
        dall'engine: ``una cosa alla volta, ci arriveremo anche a
        quello''.\footnote{Il docente non la nomina; si tratta con ogni
        probabilità della funzione \emph{High Resolution Screenshot}
        dell'editor, accessibile dal menu della viewport.}
\end{itemize}

\paragraph{Lo spadone.} Il lavoro è ``un po' meno dettagliato'', ma Mattia ha
avuto parecchi problemi con il computer. Nella stanza c'è uno spadone: ``questa
è la spada di Cloud?'' No: è uno spadone che Mattia ha davvero in camera, uno
spadone d'osso di \emph{Monster Hunter}\footnote{Nella trascrizione ``lo
spadone d'osso dimostrante''; ``dimostrante'' è quasi certamente una
storpiatura di ``di \emph{Monster Hunter}'', serie di cui Mattia aveva parlato
nel giro di presentazioni (sezione~\ref{gd8:sec:presentazioni}), e che ha
effettivamente ``tantissimi'' spadoni.} (``ottimo, fantastico''; ``ce ne sono
tantissimi, sinceramente'').

\paragraph{Daniele.} Daniele oggi non c'è: è lui che aveva avvisato. Il docente
dà comunque un'occhiata al suo lavoro: ``ottimo lavoro''.

\subsection{Francesco: la stanza senza soffitto, e il verosimile}

``Ci sono un sacco di Franceschi.'' La stanza di Francesco Muratore è ``molto
innovativa nell'avere una stanza senza soffitto'', e ``molto, molto
orizzontale''. Lo scopo dell'esercizio era fare pratica, quindi va bene; ma il
docente dà un consiglio per il futuro: fare \textbf{ambienti un po' più
credibili}. ``Il tuo letto non è così stretto'': così com'è, ``sembra più una
panchina che un letto''. Per il resto ci sono le forme, le sedie, uno zaino.

\paragraph{Non realistico: verosimile.} Serve ``un po' più di attenzione per
quanto riguarda gli spazi'', e un approccio\ldots{} ``realistico non è il
termine corretto''. Il termine che il docente userà spesso durante l'anno è
\textbf{verosimile}:
\begin{quote}
I videogiochi puntano non al realismo, ma alla \textbf{verosimiglianza}.
\end{quote}
Si torna così all'\textbf{organizzazione degli spazi}: un letto deve avere le
proporzioni di un letto perché il giocatore lo legga come tale, anche se lo
stile del gioco non è affatto realistico.

\subsection{Alberto: lo showroom dell'IKEA}

Anche Alberto ha lavorato con i \textbf{materiali}, per la maggior parte degli
oggetti (il docente si scusa: ``ci metterò un bel po' a ricordarmi i nomi di
tutti''). C'è il condizionatore. Ma la stanza è ``il vuoto totale'': nessun
poster, niente sugli scaffali vuoti. Va bene, ma la battuta è inevitabile:
``Alberto vive in uno showroom dell'IKEA''. È il rovescio del ``pizzico di
caos'' di Federica: uno spazio senza segni di vita non sembra abitato.

\subsection{Valerio: super preciso, e il problema della scala}

``Valerio invece è andato super preciso.'' Valerio racconta di essersi accorto,
provando a ``entrare'' nella stanza con il personaggio, che \textbf{era tutto
enorme}: aveva capito che i valori erano in centimetri, ma inserendo, per
esempio, 5 nel campo della scala, l'oggetto diventava di cinque metri.
Ricordiamo infatti (sezione~\ref{gd8:sec:transform}) che la \emph{posizione} è
in centimetri, ma la \emph{scala} è un \textbf{moltiplicatore} della dimensione
originale dell'oggetto: un cubo di base è largo un metro, e con scala 5 diventa
largo cinque metri (si veda anche la sezione~\ref{gd11:sec:snapping}).

``Quello non è un problema'': lo scopo era metterci mano, non preoccuparsi
troppo delle dimensioni. Il lavoro è ``un gran bel lavoro'': super dettagliato,
super preciso, con le sedie, il cestino, la \emph{tower} del computer, e una
lampada come ``tocco di classe''.\footnote{Dopo il lavoro di Valerio la
trascrizione riporta ancora un ``congratulazioni, anche tu super dettagliata'',
senza nome: potrebbe riferirsi allo stesso lavoro o all'ultimo in ordine
alfabetico. Le presenze a fine lezione confermano comunque che tutti i presenti
avevano consegnato.}

\paragraph{Finché avete energia e tempo.} ``Ci sono state un paio di persone
che hanno preso questo esercizio molto sul serio. Mi fa piacere.'' Il
consiglio: finché si hanno l'energia e il tempo per farlo, \textbf{fatelo}.
``Più andremo avanti nell'anno, meno tempo avrete. Ma è inevitabile.''
Il commento finale sulla revisione: ``ottimo, ottimo, ottimo. Soddisfazione.''

% ---------------------------------------------------------------------
\section{La lezione di oggi, e lo Starter Content che non c'è}
\label{gd11:sec:intro}

\paragraph{Teoria formale, e un po' di pratica.} Oggi si coprono concetti che
sono legati all'uso di Unreal, ma che sono anche concetti di base dello
sviluppo di videogiochi e del level design. Ci sarà anche un momento in cui si
comincia a lavorare in Unreal con lo \textbf{Starter Content}.

\paragraph{Unreal 5.8 e lo Starter Content.} Qui il docente ``incrocia le
dita'': il computer dell'aula è appena stato cambiato, e lui non ha ancora
usato la versione 5.8, perché i progetti a cui lavora usano Unity o versioni
diverse di Unreal. E nella 5.8 sono cambiate un po' di cose, tra cui una che ha
verificato proprio quella settimana: \textbf{lo Starter Content non è più
incluso con Unreal}.\footnote{Nel capitolo~\ref{gd8:cap:leveldesign}
(sezione~\ref{gd8:sec:content}) lo si poteva ancora aggiungere dal Content
Browser con \emph{Add Feature or Content Pack}. In questa lezione il docente
constata che nella 5.8 quella via non c'è più (``una volta bastava fare tasto
destro, importa starter content, hai finito. Hanno tolto sta roba, rottura
gigantesca'').} Con l'altro insegnante di level design è stato organizzato un
\textbf{pacchetto da scaricare e importare}; la procedura si vedrà più avanti
(sezione~\ref{gd11:sec:starter}).

\paragraph{Le slide sono nella cartella condivisa.} Chi preferisce seguire
sulle slide invece che sullo stream può farlo: il docente cerca di
\textbf{caricare sempre le slide in contemporanea}, quando le usa, nella
cartella condivisa. Conviene, anche perché nelle slide ci sono dei
\textbf{link} che serviranno.

% ---------------------------------------------------------------------
\section{Lavoro artistico e lavoro di gameplay}
\label{gd11:sec:arte-gameplay}

Prima della terminologia, una ``piccola introduzione'' a ``uno dei concetti più
importanti di questa disciplina'': la distinzione fra \textbf{lavoro
artistico} e \textbf{lavoro di gameplay}.

\subsection{Il termine ombrello, di nuovo}

``Cosa copre il termine level design? Quando dico level design, cosa intendo?''
Qualcuno risponde con la parola magica, ``dipende'' (``è la risposta a
praticamente ogni domanda che faccio, ma non sempre'');
la risposta giusta, vista la lezione precedente, è che level design è un
\textbf{termine ombrello} (figura~\ref{gd8:fig:ombrello}): si usa per riferirsi a
un sacco di cose diverse. Ma queste cose, ``a livello di massima'', si possono
sintetizzare in \textbf{due campi}: il campo \textbf{artistico} e il campo
\textbf{di design} (figura~\ref{gd11:fig:due-campi}).

\begin{description}
  \item[Il campo artistico.] Un level designer può trovarsi a fare
        \textbf{illuminazione}; \textbf{enrichment}, che si vedrà durante
        l'anno, cioè \textbf{arricchire gli spazi}: mettere cose sugli
        scaffali, piante e così via; in certi casi il
        \textbf{post-processing}; in certi casi gli \textbf{effetti visivi}, a
        seconda di ciò che serve: mettere fumo, fuochi in giro per una stanza.
        Sono tutte cose legate al settore artistico.
  \item[Il campo del design.] Dove metto i \textbf{nemici}? L'\textbf{encounter
        design}.\footnote{Nella trascrizione ``il counter design'': dal
        contesto (``dove metto i nemici?'') si tratta quasi certamente
        dell'\emph{encounter design}, la progettazione degli scontri.} Dove
        metto i \textbf{blocchi} che serviranno a guidare gli artisti su come
        devono essere fatte le pareti: i \textbf{blockout}, eccetera. Sono
        cose di design perché si occupano, ``non necessariamente puramente, ma
        principalmente'', dell'\textbf{esperienza del giocare al gioco}.
\end{description}

\begin{figure}[htbp]
\centering
\begin{tikzpicture}[
  font=\small,
  gd11core/.style={draw, very thick, ellipse, align=center, fill=blue!10,
                   inner sep=4pt},
  gd11art/.style={draw, thick, rounded corners=4pt, align=left,
                  text width=4.6cm, inner sep=5pt, fill=orange!12},
  gd11des/.style={draw, thick, rounded corners=4pt, align=left,
                  text width=4.6cm, inner sep=5pt, fill=green!10},
  gd11arr/.style={-{Stealth[length=2.5mm]}, thick}
]
  \node[gd11core] (ld) {\textbf{Level design}\\{\footnotesize termine ombrello}};
  \node[gd11art, below left=1.0cm and 0.3cm of ld] (art)
    {\textbf{Campo artistico}\\[2pt]
     {\footnotesize $\bullet$ illuminazione\\
      $\bullet$ \emph{enrichment}: arricchire gli spazi\\
      \phantom{$\bullet$} (scaffali, piante\ldots)\\
      $\bullet$ post-processing\\
      $\bullet$ effetti visivi (fumo, fuochi)}};
  \node[gd11des, below right=1.0cm and 0.3cm of ld] (des)
    {\textbf{Campo del design}\\[2pt]
     {\footnotesize $\bullet$ dove metto i nemici (\emph{encounter design})\\
      $\bullet$ i blockout che guidano gli artisti\\
      $\bullet$ forme e funzioni degli spazi\\
      $\bullet$ l'esperienza del giocare}};
  \draw[gd11arr] (ld) -- (art.north);
  \draw[gd11arr] (ld) -- (des.north);
  \node[font=\footnotesize\itshape, align=center, below=0.25cm of art]
    {primo anno del corso};
  \node[font=\footnotesize\itshape, align=center, below=0.25cm of des]
    {da fine primo anno;\\specialità del docente};
\end{tikzpicture}
\caption{I due macro-campi del level design. Negli studi grandi tendono a
essere coperti da persone diverse; negli studi piccoli lo stesso level
designer fa entrambe le cose.}
\label{gd11:fig:due-campi}
\end{figure}

\subsection{La specializzazione}

Ecco perché il docente fa questa precisazione: nel mondo del lavoro \textbf{i
level designer tendono a specializzarsi} in uno dei due campi.
\begin{itemize}
  \item ``Dipende, ovviamente'': in uno \textbf{studio molto piccolo} bisogna
        fare entrambe le cose, ``e a me capita di dover fare entrambi perché
        c'è poco da fare''.
  \item In uno \textbf{studio più grosso} c'è una certa specializzazione, e
        certi ambiti del lato artistico (e del level design in generale) sono
        spesso coperti \textbf{direttamente dagli artisti}, soprattutto cose
        come l'\textbf{illuminazione}: se lo studio è abbastanza grande, c'è una
        persona esperta di illuminazione che si occupa solo di quello, perché
        ad alti livelli è una cosa piuttosto complessa.
\end{itemize}

\paragraph{Il test di assunzione.} La conseguenza pratica: se un giorno si
risponde a un annuncio per \textbf{level designer} e viene chiesto di fare un
\textbf{test}, spesso si può \textbf{scegliere} se fare un \textbf{test
artistico} o un \textbf{test di gameplay}, a seconda di ciò in cui si è
specializzati.

\paragraph{Il docente, e il collega di Roma.} Il docente è specializzato nel
\textbf{design}, nel lato gameplay del level design: ``che non vuol dire che non
so niente del lato artistico, al contrario''. Il suo collega di Roma, Andrea,
proviene invece dal \textbf{lato artistico}, ed è esperto nel rendere gli spazi
e le stanze.\footnote{Nella trascrizione ``esperto nella redditività,
stanze''; il senso è quello di un esperto del lato visivo degli ambienti (la
resa, il \emph{rendering} delle stanze), ma la parola esatta non è
ricostruibile.}

\paragraph{Che cosa vuol dire per il corso.} Il \textbf{primo anno} è più
incentrato sul \textbf{lato artistico}: comporre cose in engine, imparare come
funziona l'engine, che cosa sono i materiali, l'illuminazione, eccetera. Verso
la \textbf{fine del primo anno} si comincerà a toccare il \textbf{lato del
design}. Del lato artistico il docente dice ``non è il mio punto di forza, ma
sono comunque cose che so fare''; del lato design, ``questo è il mio punto di
forza'': ``non vedo l'ora di arrivarci'', perché è ciò di cui può parlare con
più fiducia, confidenza e aneddoti. Ma intanto, ``voglio che vi ricordiate
questa distinzione principale''.

% ---------------------------------------------------------------------
\section{Il blockout}
\label{gd11:sec:blockout}

\slidefig{img/GD/gd11-blockout}{Il blockout.}{slide \emph{Intro Level Design Unreal} (2)}{fig:gd11-blockout}


``Si parte dai termini.'' Il primo è uno che gli studenti hanno già usato per il
task: il \textbf{blockout}.

\subsection{Un concetto, molti nomi}

A livello tassonomico, si possono incontrare molti nomi per la stessa cosa:
\textbf{block out}, \textbf{grey boxing}, \textbf{block mesh}, \textbf{white
boxing}. ``Questo termine ha un nome diverso praticamente in quasi ogni
studio.'' Nell'ambiente in cui lavora il docente si tende a dire
\emph{blockout}; in studi in cui ha lavorato si diceva \emph{white boxing}, in
altri \emph{grey boxing}.

\paragraph{Sono sinonimi.} Il docente torna sul punto più avanti, a proposito
di \emph{Skyrim} (sezione~\ref{gd11:sec:modulari}): capiterà di incontrare
qualcuno che dice ``no, in realtà il grey boxing è una fase del blockout che
implica questa cosa, il white boxing indica un'altra fase''. ``Io vi posso
assicurare per esperienza personale che sono usati in maniera
\textbf{completamente interscambiabile}'': \emph{white box}, \emph{white
blocking}, \emph{grey blocking}, \emph{grey box}, \emph{blockout} sono ``sempre
la stessa cosa, a livello di massima''. Saperlo evita di confondersi.

\subsection{Perché non partire dall'arte finale?}

Sulla slide c'è uno screenshot di un livello di \emph{God of War} durante lo
sviluppo:\footnote{Il docente lo indica come ``quello del 2024''. Non esiste
un \emph{God of War} uscito nel 2024: si tratta probabilmente di \emph{God of
War Ragnarök} (2022, arrivato su PC nel 2024) o del \emph{God of War} del
2018; dalla trascrizione non è possibile stabilirlo.} tutto forme semplici e
grigie. ``Perché si fanno così i livelli? Perché non partire direttamente con
la capanna e tutto quanto? Con il molo fatto bene, con la gru, con tutte le
cose pronte degli artisti?'' La risposta della classe: \textbf{prima si parte
con una bozza}.

\paragraph{La matita prima dell'inchiostro.} Il docente chiede chi sia
appassionato di fumetti e disegno, e qual è il suo fumettista preferito. Il
punto: quando si fa un disegno, di solito non si prendono subito gli
inchiostri, perché se si sbaglia qualcosa ``è tutto da buttare via''. Si parte
con la \textbf{matita}: uno sketch, si cancella, si riprova, si ricancella, si
riprova, e solo quando le matite sono finite si passa all'inchiostro. ``Il level
design funziona in una maniera molto simile.''

Quando si comincia un livello, non si fa praticamente mai così: prendere gli
asset (le case, le strade, tutto), metterli nel livello, testarlo e, se qualcosa
non funziona, dover cambiare un sacco di roba. ``No. Il livello di base è un
blockout'': \textbf{forme semplici pensate per buttare giù l'idea del livello il
prima possibile}.

\paragraph{Arrivare presto a qualcosa di giocabile.} Il motivo principale, come
suggerisce uno studente (Alessandro), è che uno degli scopi dello sviluppo di un
videogioco è \textbf{arrivare il prima possibile a qualcosa di giocabile}. Se si
aspetta che gli artisti finiscano tutto (l'erba dappertutto, le case, il cielo,
le rocce piazzate in giro) e solo a quel punto si fa il primo test, e ci si
accorge che qualcosa non funziona, per esempio che gli spazi sono troppo stretti
per il giocatore, bisogna \textbf{rifare tutto da capo}. Invece: prima si
definiscono le forme generali più ampie, e poi si gioca. ``Fondamentalmente si
fa quasi tutto il gioco così.''

\paragraph{Le realtà di produzione.} Ovviamente le realtà di produzione sono
più complesse: una \textbf{demo} o un \textbf{trailer} vanno fatti, e non si
possono fare con un blockout. ``Che cosa facciamo vedere alla gente del gioco
che stiamo sviluppando, mentre lo sviluppiamo?'' è un problema classico. Ma lo
scopo principale resta arrivare a qualcosa di giocabile il più velocemente
possibile. La questione dei blockout verrà ripresa più avanti; questa è
un'introduzione a ``una delle parti più importanti dello sviluppo di un
livello''.

\subsection{Testare idee e rispondere a domande}

Lo scopo del blockout è \textbf{testare idee e rispondere a domande}. L'esempio
del docente:
\begin{quote}
Da questo punto riesco a vedere la torre in cui devo andare? Ci metto un blocco
che ha più o meno le dimensioni di una casa. No, non si vede bene, è un po'
coperta da questa casa. Ok, la sposto più a sinistra. Perfetto, così si vede
benissimo. Andiamo avanti, e poi la sostituiremo con i lavori degli artisti,
dei modellatori 3D.
\end{quote}
``Il blockout è il vostro sketch.'' Tutti gli studenti, con il task, hanno fatto
lo sketch della propria cameretta. Lo scopo è rispondere alle possibili domande
sui problemi che vengono fuori durante la creazione di un livello, e tirare giù
idee. Ci sono anche altre fasi nella creazione di un livello, ma per ora è
questa che interessa.

\paragraph{Approfondimento: il blockout secondo Totten.} In \emph{An
Architectural Approach to Level Design}, Christopher Totten indica come errore
comune proprio quello di \textbf{presumere che il gameplay non si possa testare
finché non c'è l'arte} nell'engine: è molto più efficiente progettare gli spazi
del livello (e gli script necessari) mentre l'arte viene sviluppata, per
verificare se il progetto su carta produce davvero l'esperienza voluta. I
prototipi si costruiscono con i \textbf{primitivi dell'engine}, e in assenza di
primitivi (come in molti engine 2D) con \emph{placeholder art}, la cosiddetta
\emph{programmer art}: quadrati, cerchi, triangoli. Nel libro, l'art director
Greg Grimsby (\emph{Warhammer Online}) descrive il proprio processo in sette
passi: obiettivi di gameplay, tema visivo, punti di interesse, flusso, mappa,
\textbf{prototipo in whitebox} (``l'obiettivo è avere il guscio di base per
testare: il flusso, le visuali, i tempi di percorrenza, le linee di vista''),
e solo alla fine, una volta validato il layout con i test, l'arte finale
(``renderlo carino''). Ed è Grimsby stesso a ricordare il limite del whitebox:
un livello può funzionare anche così, ma \emph{Dead Space} ``non fa altrettanta
paura senza le luci e l'atmosfera'' dell'arte: il blockout risponde alle
domande di gameplay, non a quelle di atmosfera.

% ---------------------------------------------------------------------
\section{Due esempi: Call of Duty e Uncharted}
\label{gd11:sec:esempi-blockout}

\slidefig{img/GD/gd11-esempi}{Blockout e livello finito: Call of Duty e Uncharted.}{slide \emph{Intro Level Design Unreal} (2)}{fig:gd11-esempi}


``A me piace cercare di portare esempi pratici, quando riesco a scovarli fra
archivi, dichiarazioni, o chiedendo ad amici o persone che hanno conoscenze.''

\subsection{Una mappa multiplayer di Call of Duty}

Il primo esempio è una mappa multiplayer di \emph{Call of Duty} (``quello del
2019'', cioè \emph{Call of Duty: Modern Warfare}), realizzata da un level
designer di cui il docente mostra il lavoro.\footnote{Nella trascrizione la
mappa è chiamata ``DOCS'' e il level designer ``Brian Baker''; né l'una né
l'altro sono verificabili dalla trascrizione, e li si riporta solo come
indicazione. Il docente aggiunge: ``non chiedetemi quale studio, perché a
\emph{Call of Duty} ci lavorano tipo trenta studi diversi''.} In alto c'è la
versione in blockout, in basso la versione finale (figura~\ref{gd11:fig:cod}).
``Che cosa notate?''

\paragraph{Materiali, griglia, colori.} Anzitutto, nel blockout c'è una
\textbf{scelta di materiali precisa}:
\begin{itemize}
  \item una \textbf{griglia}, che aiuta a tenere presenti le distanze, la
        \textbf{metratura} di ciò che si sta facendo;
  \item dei \textbf{colori} pensati per \textbf{segnalare} qualcosa in
        particolare.
\end{itemize}
A che cosa servano esattamente quei colori ``bisognerebbe chiederlo
direttamente al level designer che li ha scelti''. Secondo il docente, lo scopo
principale è aiutare a notare, in un ambiente in cui tutto ha la stessa
texture, che lì c'è qualcosa: per esempio \textbf{una finestra}. E la cosa è
stata rispettata: la finestra del blockout e quella della versione finale sono
lo stesso edificio. Nel blockout interessano solo le forme; il lato estetico
``assolutamente zero''.

\paragraph{Perché segnalare una finestra?} L'orlo \textbf{rosso} intorno
all'apertura serve a dire al giocatore: \emph{occhio, lì c'è una finestra}. Ma
perché un level designer vorrebbe farlo notare? La finestra è al primo piano di
un edificio, e in \emph{Call of Duty} del 2019 il giocatore non ci può
arrivare.\footnote{Uno studente obietta che in un altro capitolo della serie il
personaggio ha il doppio salto (la trascrizione dice ``Wolverine Warfare'',
probabilmente uno dei capitoli con movimento avanzato, come \emph{Advanced
Warfare} o \emph{Infinite Warfare}); il docente risponde che ``è un discorso a
parte''.} Lo scopo, in realtà, è segnalare: \textbf{da lì potrebbe spararti
qualcuno}. ``Tieni presente che lì magari c'è qualcuno appostato che ti sta
puntando col fucile.'' E questa cosa è stata \textbf{mantenuta nella versione
finale}, perché una delle cose che fa un blockout è \textbf{segnalare
l'intento} delle varie parti.

\paragraph{La luce come segnale.} Nella versione finale, guardando
l'\textbf{illuminazione}, si nota un'attenzione particolare: il giocatore che
arriva correndo e svolta l'angolo deve notare che lì c'è un colore diverso
rispetto alla parete. L'hanno ottenuto inserendo una \textbf{luce}
(probabilmente dentro c'è una lampadina) che fa notare: ehi, lì c'è
un'apertura nel muro, potrebbe esserci qualcuno che ti aspetta dietro l'angolo
per spararti. ``Questo è il lavoro del level design: fare attenzione a questo
tipo di dettagli.''

\begin{figure}[htbp]
\centering
\begin{tikzpicture}[
  font=\small,
  gd11wall/.style={draw, thick, fill=gray!25},
  gd11win/.style={draw=red!80!black, line width=1.6pt, fill=black!70},
  gd11winc/.style={draw, thick, fill=blue!15},
  gd11lab/.style={font=\footnotesize, align=left},
  gd11arr/.style={-{Stealth[length=2.5mm]}, thick}
]
  % Blockout
  \node[font=\small\bfseries] at (1.6,3.3) {Blockout};
  \draw[gd11wall] (0,0) rectangle (3.2,2.8);
  \foreach \x in {0.4,0.8,...,2.8} \draw[gray!55, thin] (\x,0) -- (\x,2.8);
  \foreach \y in {0.4,0.8,...,2.4} \draw[gray!55, thin] (0,\y) -- (3.2,\y);
  \draw[gd11win] (1.2,1.6) rectangle (2.0,2.3);
  \node[gd11lab, anchor=north] at (1.6,-0.15)
    {un solo \textbf{buco}: la funzione\\(``da qui ti tengono d'occhio'')};
  % Freccia
  \draw[gd11arr] (3.6,1.4) -- node[above, font=\footnotesize]{test,}
                              node[below, font=\footnotesize]{arte} (5.0,1.4);
  % Finale
  \node[font=\small\bfseries] at (6.9,3.3) {Versione finale};
  \draw[thick, fill=brown!20] (5.3,0) rectangle (8.5,2.8);
  \draw[gd11winc] (5.6,1.6) rectangle (6.3,2.3);
  \draw[gd11win, draw=black, line width=0.8pt, fill=black!75] (6.55,1.6) rectangle (7.25,2.3);
  \draw[gd11winc] (7.5,1.6) rectangle (8.2,2.3);
  \fill[yellow!80!orange, opacity=0.55] (6.9,1.95) circle (0.22);
  \node[gd11lab, anchor=north] at (6.9,-0.15)
    {la stessa apertura, tra due finestre\\\textbf{chiuse} e non sfondabili; una\\\textbf{luce} interna la fa notare};
\end{tikzpicture}
\caption{Ricostruzione schematica dell'esempio di \emph{Call of Duty}
(2019): nel blockout l'apertura è solo un buco segnalato in rosso, che dice
``da lì potrebbero spararti''; nella versione finale l'intento è mantenuto,
l'apertura diventa una finestra resa verosimile da altre finestre chiuse, e una
luce la rende leggibile a chi svolta l'angolo.}
\label{gd11:fig:cod}
\end{figure}

\paragraph{Funzione prima, estetica poi.} C'è anche una differenza tra
blockout e versione finale: nel blockout c'è \textbf{una sola finestra aperta},
mentre nella versione finale ci sono più finestre. Il docente si spinge oltre:
quando si faceva questo livello, nella fase delle forme, probabilmente non
c'era neanche l'idea di una finestra. L'idea era: \textbf{questo è il buco da
cui puoi tenere d'occhio il cortile}. ``Quello è puro gameplay, quella è pura
considerazione di gameplay.'' Successivamente si sarà deciso che quel buco era
una finestra, e per radicarlo in un mondo \textbf{verosimile} sono state
aggiunte altre due finestre, chiuse, che però non si possono sfondare. A
livello funzionale è un buco; a livello estetico, una serie di finestre.

\paragraph{Perché il blockout costa meno.} ``Il blockout ha uno scopo
puramente funzionale.'' Lo si fa, lo si testa, si gioca la mappa un paio d'ore
con altre persone (``ovviamente tutto quello che sto dicendo è
ipersemplificato: a \emph{Call of Duty} avranno procedure molto più complesse'').
L'apertura in quel punto non va bene? La si sposta un po' più a destra, si
sistemano altre cose nel blockout, si va avanti. Se invece si fosse partiti
dall'\textbf{edificio finito}, spostare quella finestra un po' più a destra
``sarebbe stato un enorme casino'', perché tutto il resto del modello andrebbe
cambiato. Prima il blockout; poi, in base a ciò che viene fuori dai test e dai
\emph{playthrough}, si dice ``ok, questo è il livello, non tocchiamo quasi più
niente''. Solo a quel punto si fanno i \textbf{modelli}: in base allo spazio da
rispettare, alle forme stabilite, alle \textbf{funzioni legate al gameplay}
date dal blockout iniziale del level designer.

\subsection{Il level designer dà ordini agli artisti?}

``Una domanda sorge spontanea, vi leggo nel cervello: ma Stefano, questo vuol
dire che un level designer dà gli ordini agli artisti? Si metteranno a fare la
casa in base a quello che decido io col livello?'' La risposta è: \textbf{dipende}.
Dalla struttura del team, dalle funzioni, dal fatto che la casa sia un
\textbf{modello unico} o un insieme creato con \textbf{unità modulari} che si
possono riarrangiare ``tipo Lego'' (sezione~\ref{gd11:sec:modulari}). Ci sono
varie procedure. Ma il messaggio è un altro:
\begin{quote}
Prima di tutto si decide il \textbf{gameplay}; prima di tutto si decide la
\textbf{funzione} degli spazi; prima di tutto si decidono le \textbf{forme}
degli spazi. Dopodiché questi spazi vengono \textbf{estetizzati}.
\end{quote}

\paragraph{Il gameplay regna sovrano.} Chi si è mai chiesto come sia possibile
che ci siano videogiochi in cui \textbf{la trama viene scritta pochi mesi prima
della fine} dello sviluppo, invece di essere decisa dall'inizio, ha qui la
risposta: in questo tipo di processo \textbf{il gameplay regna sovrano}. Prima
ci si assicura di fare un gioco funzionale, efficace, entusiasmante; poi si
decide la veste da dargli.

\paragraph{Ma non a cascata.} Con una precisazione importante: la \textbf{veste
artistica} è ovviamente decisa \textbf{in precedenza}. ``Non è che prima si
fanno tutti i blockout e poi si decide: ah, questo è fantascientifico o
medievale?'' E molto spesso, per questioni di produzione, \textbf{non si lavora
quasi mai a cascata} (si fa una cosa, poi l'altra, poi l'altra): quante più
cose possibile procedono \textbf{in parallelo}, ``perché il tempo è denaro''.

\subsection{Uncharted: dichiaro che è una cascata}

Il secondo esempio è l'aspetto che aveva \emph{Uncharted} durante lo sviluppo.
Il level designer dice: ``qui ci voglio mettere una piccola cascata. Questa è
una cascata. \textbf{Dichiaro che è una cascata}.'' E poi sta agli artisti, al
level designer che si occupa del lato artistico o allo stesso level designer
andare a prendere gli asset, piazzarli lì e fare in modo che siano fatti bene.

``Voglio che vi fissiate in testa che \textbf{questo è l'aspetto che hanno tutti
i videogiochi nel corso del loro sviluppo}.'' L'argomento tornerà, con lezioni
dedicate.

% ---------------------------------------------------------------------
\section{I cinque strumenti del blockout}
\label{gd11:sec:strumenti}

\slidefig{img/GD/gd11-strumenti}{I cinque strumenti del blockout.}{slide \emph{Intro Level Design Unreal} (2)}{fig:gd11-strumenti}
\slidefig{img/GD/gd11-strumenti2}{Gli strumenti del blockout in dettaglio: primitive, kit modulari, brush e scultura.}{slide \emph{Intro Level Design Unreal} (2)}{fig:gd11-strumenti2}


``Restando in argomento blockout, come si creano i blockout?'' Esistono
\textbf{cinque strumenti principali} (tabella~\ref{gd11:tab:strumenti}). Prima
di vederli uno per uno, il messaggio che il docente vuole far passare:
\begin{quote}
Questi strumenti \textbf{non sono mutualmente esclusivi}. Possono essere usati
tutti insieme, a seconda di quello che dobbiamo fare. Il trucco è prendere
\textbf{lo strumento più veloce per lo scopo}.
\end{quote}
E non servono solo al blockout: molti di questi strumenti vengono usati anche
nelle \textbf{fasi finali}, per costruire i livelli veri e propri e
posizionare gli oggetti.

\begin{table}[htbp]
\centering
\small
\begin{tabular}{@{}p{2.4cm}p{3.6cm}p{3.6cm}p{3.6cm}@{}}
\toprule
\textbf{Strumento} & \textbf{Che cos'è} & \textbf{Vantaggi} & \textbf{Svantaggi} \\
\midrule
Primitivi & forme di base dell'engine (cubo, cilindro, sfera, cono, piano, scale\ldots) &
  semplicissimi da usare, veloci & limitati alle forme disponibili (una catena curva richiede una vita) \\
\addlinespace
Kit modulari & set di pezzi a misura (pareti, porte, finestre, angoli, pilastri) che si incastrano &
  velocissimi; spazi già a misura del gioco; spesso già i pezzi finali &
  rigidi: difficile uscire dagli schemi (una stanza circolare è quasi impossibile) \\
\addlinespace
Brush e modellazione & modifica diretta della geometria (vertici, facce, estrusioni), nell'engine o in Blender &
  qualunque forma; con la pratica, velocità fenomenale &
  più lenti, richiedono precisione; curva di apprendimento \\
\addlinespace
Scultura (terreno) & editor di terreno: alzare, abbassare, allargare, lisciare una superficie &
  l'unico modo sensato per colline e terreni naturali &
  non lavora su facce e vertici singoli; adatto solo al terreno \\
\addlinespace
Spline & curva definita da punti, lungo cui l'engine dispone e piega una o più mesh &
  strade, cavi, recinzioni curve senza modellarli a mano; modificabili in un attimo &
  più tecniche; se la mesh si allunga troppo si vedono stiramenti \\
\bottomrule
\end{tabular}
\caption{I cinque strumenti principali per costruire un blockout (e, in buona
parte, anche i livelli finali).}
\label{gd11:tab:strumenti}
\end{table}

\subsection{I primitivi}
\label{gd11:sec:primitivi}

I primi sono quelli già usati nel task: i \textbf{primitivi}. ``Come dice il
termine, un primitivo è una forma base.'' Le forme di base di Unreal, che ``sono
comuni in praticamente tutti gli engine'', sono forme che \textbf{non possono
essere create da altre forme}: sono gli \textbf{elementi di base}. Il cubo, il
cilindro, la sfera, il cono, il piano (``semplicemente piatto, per pavimenti,
pareti e cose del genere''); Unreal offre anche le scale e altre
forme.\footnote{Il docente distingue di passaggio le \emph{forme} di base dalle
\emph{geometrie} (``le geometrie sono una cosa diversa, anche quelle prima o
poi le vedremo''): si riferisce ai \emph{brush} di geometria, visti nella
sezione~\ref{gd11:sec:brush}.}

È il modo più basilare di fare un blockout: ``hai il tuo cubo, lo stretchi;
vuoi fare la colonna, prendi sei cilindri, pam pam pam, colonne fatte, e in
cinque minuti ho fatto la prima versione del tempio in cui voglio ambientare la
mia boss fight''.

\paragraph{Vantaggi e svantaggi.} I primitivi sono \textbf{molto semplici da
usare} e da sfruttare: si va relativamente veloci. Lo svantaggio è che
\textbf{si è limitati alle forme che si hanno}. Se si sta facendo un livello
``alla \emph{Dark Souls}'', ``alla \emph{Elden Ring}'', dove c'è un punto con
delle \textbf{catene gigantesche} su cui si deve correre, fare una catena curva
con i primitivi ``ci metto una vita''. Per questo ci sono altri strumenti.

\paragraph{Approfondimento: primitivi e griglie secondo Totten.} Totten osserva
che i primitivi di un engine tipico comprendono cubi (o parallelepipedi),
sfere, cilindri e piani, a volte anche scale e coni, e che con essi si può
costruire ``una varietà sorprendente'' di geometria: in molti engine quasi
tutta la geometria spaziale più grande di un livello può essere fatta solo con
primitivi. Un vantaggio ulteriore è che i primitivi \textbf{seguono un sistema
di misure standard} dettato dall'engine. Il caso estremo è l'editor
\textbf{Hammer} di Valve (lo stesso citato dal docente nella
sezione~\ref{gd11:sec:brush}), che impone di lavorare su una griglia con lo
snapping sempre attivo, per evitare buchi nel livello che ne impedirebbero la
compilazione: gli oggetti scattano a misure in \textbf{potenze di due} (16, 32,
64, 128 unità\ldots).

\subsection{I kit modulari}
\label{gd11:sec:modulari}

Per strutture più complesse, più particolari, il secondo strumento sono i
\textbf{kit modulari}. Sulla slide c'è una serie di pezzi: ``che cosa sono?''
Pareti, e i buchi al loro interno sono \textbf{porte} e \textbf{finestre}. Con
\textbf{dodici elementi} si può costruire l'interno di un intero palazzo, cioè
un insieme di stanze, perché quei pezzi coprono \textbf{tutte le possibili
connessioni} tra una stanza e l'altra, più il pavimento: angoli, porte,
pilastri per gli angoli che si incontrano sul lato esterno, perfino spazi di
dimensioni diverse a seconda di come si vuole organizzare. Si vogliono delle
colonne al centro delle stanze? Si piazzano i pezzi giusti.

\paragraph{Lego e tasselli.} ``Se i primitivi sono Lego, i kit modulari sono le
caselle con cui si costruisce uno spazio.''

\paragraph{Vantaggi.} Sono \textbf{molto più veloci}: si costruiscono spazi in
pochissimo tempo. In più, essendo kit già sistemati (in teoria) con
\textbf{dimensioni utili per il gioco}, non ci si deve neanche preoccupare degli
spazi: automaticamente si sa che funzionano. Di kit modulari ce ne sono molti
tipi diversi; e non contengono solo forme basiche, ma anche scale e forme più
complesse, componibili velocemente. Danno una \textbf{visione del livello molto
vicina a quella finale}, perché manca davvero solo il lato estetico. E in un
gioco come \emph{Skyrim}, il kit modulare \textbf{è già ciò che si usa nella
versione finale}.

\paragraph{Svantaggi.} ``Fatemi una stanza circolare con questo kit.'' È
sostanzialmente \textbf{impossibile}: non si hanno gli strumenti. I kit modulari
permettono di fare un sacco di cose molto velocemente, ma rendono \textbf{molto
più difficile uscire dagli schemi}: ``tendono a peccare di creatività'', perché
con quegli strumenti non si può fare più di tanto. Ma, di nuovo, si combinano
gli strumenti: se c'è uno spiazzo e poi si entra in un palazzo con delle stanze,
le stanze le si fa velocemente con il kit, lo spiazzo con altri metodi.
L'importante è fare tutto il prima possibile, per poter subito testare le
meccaniche, vedere se il livello funziona, buttarci dentro i nemici e provare.

\paragraph{Sistemi e contenuti.} Per capire dove i kit modulari brillano, il
docente richiama una distinzione: i \textbf{sistemi} sono le meccaniche di
gameplay; i \textbf{contenuti} sono ``le cose a cui effettivamente giochi'': i
livelli, le armi, tutto ciò che ``fa parte, lo dice il termine, del contenuto
del gioco stesso''. Giochi come \emph{Skyrim} o \emph{Elden Ring} richiedono una
\textbf{quantità enorme di contenuti}, ``veramente vastissima'', e ci sono punti
in cui bisogna essere intelligenti nello sfruttare ciò che si ha a disposizione.

\paragraph{I dungeon di Elden Ring.} ``Chi ha giocato a \emph{Elden Ring}? Qual
è la parte più ripetitiva?'' Uno studente cita una zona della capitale; il
docente punta su altro: i \textbf{dungeon segreti} (catacombe, grotte), ``tutti
quasi uguali''. ``Posso quasi garantire che sono stati fatti tutti con i kit
modulari.''

\paragraph{I roguelike.} Un passaggio ulteriore: i \textbf{roguelike} che
costruiscono le stanze in maniera casuale usano proprio questo approccio:
\textbf{pezzi di stanza riorganizzati} da un \textbf{algoritmo} che li genera
ogni volta.

\paragraph{Il kit di Skyrim.} Sulla slide c'è il kit modulare dei dungeon di
\emph{Skyrim}: contiene sia parti \textbf{decorative} sia parti più
\textbf{pratiche}, tutte le pareti usate per costruire i dungeon. Per chi vuole
approfondire, nella slide c'è il link alla \textbf{trascrizione di un talk}
tenuto nel 2013 da uno dei level designer di \emph{Skyrim}, che spiega come sono
stati realizzati i dungeon con queste metodologie (dimensioni, studi, eccetera)
e come è stato creato il kit.\footnote{Il docente non dice il nome. Si tratta
quasi certamente del talk ``\emph{Skyrim's Modular Approach to Level Design}''
tenuto alla GDC 2013 da Joel Burgess (Bethesda), la cui trascrizione è
pubblicata sul blog dell'autore.} Loro lo chiamano \emph{grey box}: ``giusto per
tornare all'argomento che ognuno chiama il blockout con il termine che
preferisce''.

\paragraph{Approfondimento: la modularità secondo Totten.} Totten ricorda che la
modularità viene dall'architettura: dalla Rivoluzione industriale gli edifici si
costruiscono con \textbf{elementi prefabbricati}, assemblati in cantiere; Mies
van der Rohe e Philip Johnson progettavano con componenti prefabbricati di
acciaio e vetro, e Le Corbusier ha progettato l'Unité d'Habitation con
appartamenti come unità modulari impilate. Totten usa proprio la metafora del
Lego: con un set di pezzi modulari, i livelli si assemblano ``come set di Lego
invece che come opere d'arte su misura difficili da cambiare''. La modularità
aiuta anche le \textbf{metriche}: nel gioco \emph{SWARM!} le piastrelle modulari
larghe due unità erano la misura minima su cui il giocatore poteva restare in
equilibrio senza che il gioco risultasse ingiusto. E, come gli architetti
tengono traccia dei pezzi prefabbricati con i \emph{building schedule} (tabelle
che indicano tipo, misure e posizione di porte, finestre, impianti), il level
designer può documentare i propri pezzi modulari allo stesso modo.

\subsection{Brush e modellazione}
\label{gd11:sec:brush}

``Ok, questa è dove le cose si fanno interessanti.'' Il terzo strumento sono i
\textbf{brush}, o la \textbf{modellazione}. La domanda che potrebbe sorgere:
se si sta facendo una cosa da zero, e l'unica cosa che interessa è la forma
generale del livello per ``trovare la quadra iniziale'', \textbf{perché non
usare direttamente uno strumento di modellazione 3D}, invece di giocare con le
forme o con i kit modulari?

La risposta è che \textbf{si può fare}. Alcuni engine permettono di
\textbf{modificare direttamente la geometria}, come se si lavorasse in ``una
versione ipersemplificata di Blender o Maya''. Per fare un letto, per esempio,
si ritaglia una sezione di un blocco, la si prende e la si tira su: ``e ho fatto
il mio letto''. Si fa un \textbf{editing diretto} di ciò che si sta costruendo
nel livello.

\paragraph{Blender e Unity.} Una delle pipeline più comuni quando si lavora in
\textbf{Unity} (``che non ci interessa, perché noi lavoreremo in Unreal, ma è
importante saperlo nel caso vi capiti'') è lavorare direttamente in
\textbf{Blender}: molto spesso i level designer costruiscono il livello in
Blender, che salva direttamente nel progetto Unity, e fanno così l'intero
blockout del livello molto velocemente.

\paragraph{Vantaggi e svantaggi.} Il vantaggio: \textbf{si può fare quello che
si vuole}, letteralmente qualunque forma. Non c'è il problema ``non ho il
primitivo giusto'': la forma la si crea, la si modifica, la si gira come si
vuole. Lo svantaggio (``non è una questione di peso, in realtà''): è \textbf{un
po' più lungo}, bisogna essere più \textbf{precisi}, bisogna \textbf{saper usare
questi strumenti}. Ma c'è un ``ma'' importantissimo: se si è disposti a
impararli, questi strumenti permettono di andare a \textbf{una velocità
fenomenale}, soprattutto per i blockout. Il problema è la \textbf{curva di
apprendimento}: Blender richiede un po' di tempo per essere imparato (nel corso
lo si toccherà); ma si vedrà soprattutto come fare questo lavoro dentro Unreal,
con i \textbf{brush di geometria}.

\paragraph{Valve e Hammer.} Chi ha giocato a un gioco di \textbf{Valve}
(\emph{Half-Life}, \emph{Half-Life 2}, \emph{Half-Life: Alyx},
\emph{Deadlock}: ``ve li consiglio tutti'')\footnote{Nella trascrizione
``AlphaFlash, AlphaFlash, AlphaFlash, Alix, Deadlock'': i primi tre sono
evidentemente i titoli della serie \emph{Half-Life}.} ha visto il risultato di
questo approccio. ``C'è un motivo se Valve fa un gioco ogni dieci anni: se lo
fa, vuol dire che sono riusciti a fare il proprio lavoro, dato che hanno soldi
infiniti.'' Valve lavora con \textbf{Hammer}, un editor proprietario per il
proprio engine (\textbf{Source}): in Hammer è \textbf{tutto editing diretto}
delle mesh e delle forme nel livello, da parte del level designer.

\paragraph{Brush e modello 3D.} Qual è la differenza tra un \textbf{brush} e un
\textbf{modello 3D}? ``Lo vedremo'': sono due sistemi diversi. Per ora basta
sapere che \textbf{è possibile creare forme direttamente nell'engine}, ed è un
altro dei modi di fare il blockout.

\subsection{La scultura del terreno}
\label{gd11:sec:scultura}

Così come si può modificare la geometria tirando punti, vertici, facce,
estrudendo e suddividendo, \textbf{si può anche scolpire}. Il quarto strumento è
l'\textbf{editor di terreno} di Unreal (il \emph{Landscape}), che si vedrà
durante l'anno.

Se si sta facendo un livello ambientato su una \textbf{collina}, la collina non
la si può fare con dei cubi: non si riesce a fare un terreno scosceso con i
primitivi. Per i paesini sulla collina di \emph{Elden Ring} si mettono cubi
dappertutto? ``Perdo un sacco di tempo e il risultato è orrendo.'' No: si usa un
editor che permette di \textbf{alzare, modificare, allargare, scolpire} una
superficie.

\paragraph{Scolpire non è modellare.} Scolpire significa che non si stanno
tirando facce o vertici individuali, come con i brush o la modellazione: ``passo
col mouse, dove passo col mouse si alza un pochino, poi lo liscio un po'
dall'altra parte''. È così che si fanno gli \textbf{ambienti naturali}.

\subsection{Le spline}
\label{gd11:sec:spline}

``Ultimo ma non ultimo, e molto importante'': le \textbf{spline}. È uno
strumento un po' più tecnico, ma è il quinto dei cinque, e serve sia alla
concezione iniziale del livello sia al \textbf{posizionamento degli oggetti}
nelle fasi finali.

\paragraph{Il problema della strada.} Sulla slide c'è una strada. Se la strada è
dritta è facile: un pezzo di strada, poi un altro pezzo. Ma se fa una curva?
Si chiede agli artisti di modellare a mano ogni possibile strada del gioco? ``A
livello pratico magari è fattibile.'' Ma se poi un'area viene cancellata, o la
strada deve curvare a sinistra invece che a destra, bisogna rifare tutto da
capo. ``La risposta, ovviamente, è no.'' E la risposta ``è nell'intrinseco della
domanda'': come si gestiscono le cose che devono essere \textbf{piegate} dentro
un videogioco?

\paragraph{Che cos'è una spline.} Attenzione: ``la spline qui non è la strada,
\textbf{la spline è la linea}''. Una spline è \textbf{una curva definita da dei
punti}, che curva in base ai dati contenuti in quei punti: la
\textbf{rotazione}, le dimensioni dei \textbf{vettori di ingresso e di uscita}
(le tangenti), e altro. A questa curva vengono poi assegnate \textbf{una o più
mesh}, e l'engine fa in modo che la mesh \textbf{segua la curva}
(figura~\ref{gd11:fig:spline}). Che cosa sia esattamente una mesh si vedrà tra
poco (sezione~\ref{gd11:sec:mesh}); per ora, ``pensate a un modello 3D''.

\paragraph{Un pattern?} Uno studente chiede se si tratti di un
\textbf{pattern}, cioè sempre lo stesso elemento copiato. ``Molto spesso sì'', e
il docente ringrazia, perché stava per arrivarci: molto spesso la spline usa
\textbf{più oggetti ripetuti}, attaccati ai bordi uno dell'altro e modificati in
modo da piegarsi.

\paragraph{Il problema dello stretch.} Il perché lo spiega alla lavagna. Si
prenda una superficie divisa in poligoni. Se la si \textbf{allunga} (la si
``stretcha''), i dati del modello 3D non cambiano: si sta solo dicendo
``allungalo''. Quindi gli spazi interni dei poligoni vengono stirati anch'essi:
\textbf{non si aggiungono dettagli}. Se una spline avesse un solo modello da
deformare, si noterebbe che è ``superstretchato'': le texture stirate, i
poligoni sempre più grandi. Per questo la spline \textbf{attacca tra loro più
copie} del modello, e in curva le \textbf{comprime leggermente}. Le facce
restano le stesse, ma se i modelli sono fatti bene e con cognizione di causa
(abbastanza densi di poligoni) e la modifica è minima, ``il giocatore non lo
nota''. Poi, certo, guardando con attenzione in qualunque videogioco si notano
le giunture da qualche parte.

\paragraph{I cavi.} Lo stesso vale per i \textbf{cavi}: non esiste che, per ogni
possibile modo di posare un cavo nel gioco, il modellatore 3D faccia un cavo
diverso. Oltre al lavoro, questo approccio farebbe diventare le
\textbf{dimensioni del gioco} enormi. Si fa \textbf{un cavo dritto}, e poi con
metodi come le spline (``non per forza le spline, ma la spline è il metodo di
base più utilizzato'') se ne dispongono e piegano varie copie.

\begin{figure}[htbp]
\centering
\begin{tikzpicture}[
  font=\small,
  gd11pt/.style={circle, fill=blue!70!black, inner sep=1.8pt},
  gd11tan/.style={-{Stealth[length=2mm]}, thick, blue!50},
  gd11lab/.style={font=\footnotesize, align=center}
]
  % (a) un solo modello stirato
  \node[font=\small\bfseries] at (2.2,2.4) {(a) un solo modello, allungato};
  \foreach \i in {0,...,3} {
    \draw[thick, fill=gray!20] ({\i*1.1},0.9) rectangle ({\i*1.1+1.1},1.7);
  }
  \node[gd11lab] at (2.2,0.45) {i poligoni si stirano,\\la texture si deforma};
  \draw[thick, fill=gray!20] (0.9,1.9) rectangle (1.3,2.05);
  % (b) spline con piu copie
  \node[font=\small\bfseries] at (8.6,2.4) {(b) spline con più copie};
  \draw[very thick, blue!70!black] plot[smooth, tension=0.8]
       coordinates {(6.0,0.6) (7.3,1.6) (8.8,1.7) (10.2,0.9) (11.2,0.4)};
  \foreach \p/\a in {(6.0,0.6)/40, (8.8,1.7)/0, (11.2,0.4)/-25} {
    \node[gd11pt] at \p {};
  }
  \draw[gd11tan] (6.0,0.6) -- ++(40:0.9);
  \draw[gd11tan] (8.8,1.7) -- ++(0:0.9);
  \draw[gd11tan] (8.8,1.7) -- ++(180:0.9);
  \foreach \x/\y/\r in {6.35/0.95/38, 6.95/1.38/30, 7.6/1.68/10, 8.25/1.76/2,
                         8.9/1.7/-6, 9.55/1.45/-22, 10.15/1.0/-30, 10.75/0.62/-26} {
    \draw[thick, fill=gray!25, rotate around={\r:(\x,\y)}] (\x-0.3,\y-0.12) rectangle (\x+0.3,\y+0.12);
  }
  \node[gd11lab] at (8.6,-0.3) {punti di controllo con rotazione e tangenti\\(vettori di ingresso e uscita); copie della mesh\\attaccate tra loro e leggermente compresse in curva};
\end{tikzpicture}
\caption{Perché le spline usano più copie di una mesh: allungare un solo modello
(a) stira poligoni e texture senza aggiungere dettaglio; la spline (b)
dispone lungo la curva più copie di un pezzo dritto, deformandole quel poco che
basta perché il giocatore non lo noti.}
\label{gd11:fig:spline}
\end{figure}

\paragraph{Una lezione frontale.} ``Queste sono le cinque metodologie, per darvi
un'infarinatura iniziale. Tutto chiaro fino ad adesso?'' Il docente si rende
conto che è una delle lezioni più frontali che deve fare, ``ma vi devo insegnare
tutta una serie di termini e c'è poco da fare''.

% ---------------------------------------------------------------------
\section{Che cosa si farà nel primo anno}
\label{gd11:sec:primo-anno}

\paragraph{Kit modulari e primitivi.} Nel primo anno si lavorerà
\textbf{principalmente con kit modulari e primitivi}. Lo scopo è imparare il più
possibile a usare gli strumenti. Si è deciso di partire dal \textbf{lato
artistico} perché l'obiettivo è sviluppare una \textbf{conoscenza istintiva} di
ciò che si usa: ``in modo che se ti dici `voglio fare un livello così', non devi
perdere tempo a cercare come si fa: sai già come si fa, e vai''.

\paragraph{L'allenamento dell'occhio.} Quest'anno, quindi, si terrà d'occhio la
\textbf{costruzione} e l'\textbf{estetica} dei lavori. Ci sarà una forte
concentrazione su:
\begin{itemize}
  \item attenzione a \textbf{ciò che si mette insieme};
  \item attenzione alla \textbf{coerenza} del lavoro;
  \item attenzione agli \textbf{spazi}: ``parleremo tanto di suddivisione degli
        spazi'';
  \item attenzione all'\textbf{occhio}: ``questo sarà l'anno
        dell'occhio''.\footnote{Nella trascrizione ``questo sarà l'armo
        dell'occhio''; la lettura più plausibile è ``l'anno dell'occhio'',
        cioè dell'allenamento dello sguardo.}
\end{itemize}
Prima si impara a usare gli strumenti, poi li si applicherà più precisamente ai
videogiochi.

% ---------------------------------------------------------------------
\section{Asset e actor, di nuovo}
\label{gd11:sec:asset}

Si passa alla terminologia dell'editor, partendo da un termine già visto:
\textbf{asset} (sezione~\ref{gd8:sec:content}). Uno studente risponde
correttamente, e il docente precisa:
\begin{description}
  \item[Asset] un \textbf{file salvato nel progetto}, su disco;
  \item[Actor] un'\textbf{istanza} di quell'asset \textbf{creata all'interno
        del gioco}, nel livello.
\end{description}

Mentre Unreal si apre sul computer nuovo (``è totalmente vuoto'', nessun
progetto già pronto, e serve l'accesso a internet), il docente ricorda la
disposizione dell'editor vista la volta precedente:
\begin{itemize}
  \item in basso c'è il \textbf{Content Browser}: gli asset \textbf{del
        progetto}, interni al progetto;
  \item a sinistra (\emph{Place Actors}) ci sono gli \textbf{asset di sistema
        di Unreal}: sono sempre asset, ma sono quelli \textbf{dell'engine}. Le
        luci, le forme di base e così via sono già presenti in Unreal, e non c'è
        bisogno di rifarle ogni volta che si crea un progetto nuovo.
\end{itemize}
Una conseguenza, emersa da una domanda durante la pausa: gli asset
dell'engine, quelli presi dal pannello a sinistra, \textbf{non si possono
modificare a livello di asset}, ``perché sono di Unreal''; quelli modificabili
sono quelli del Content Browser.

% ---------------------------------------------------------------------
\section{Coordinate globali e coordinate locali}
\label{gd11:sec:coordinate}

Nell'attesa, il docente passa alla lavagna (Mattia conferma che da casa la si
vede) per una distinzione che ``tende a creare confusione all'inizio'': quella
tra \textbf{world coordinates} (coordinate globali) e \textbf{local
coordinates} (coordinate locali). ``So che probabilmente molti di voi già lo
sanno, ma questa lezione mi serve principalmente per assicurarmi che sulla
carta, in teoria, sappiate tutti di cosa sto parlando.''

\subsection{L'origine del mondo}

Si disegna un mondo di gioco, con un giocatore. ``Questo è il punto
$(0,0,0)$. Che cos'è il punto zero?'' Qualcuno propone il punto di spawn: no.
Il \textbf{punto di spawn} è il punto in cui il giocatore viene messo in gioco,
e non c'entra niente. Il punto $(0,0,0)$ è l'\textbf{origine delle coordinate}
del mondo di gioco: il punto da cui partono tutte le coordinate.

\paragraph{Quale asse è l'alto?} Le coordinate sono $x, y, z$, ma l'ordine e il
significato degli assi \textbf{cambiano a seconda dell'engine} (e a volte del
progetto): in Unreal \textbf{$z$ è l'altezza}; in altri engine $z$ è la
\textbf{profondità} e l'altezza è $y$. ``Lungo dibattito su quale sarebbe la
soluzione migliore; non esiste un modo condiviso.'' Bisogna farci attenzione
quando si lavora con engine diversi: ciò che conta è sapere quale asse è
altezza, quale larghezza e quale profondità.

\paragraph{Coordinate globali.} Un giocatore che si trova, ``a naso'', a
cinque unità dall'origine lungo un asse avrà coordinate $(5,0,0)$ (``pessima
visualizzazione, perché è in 2D''). In 3D, un oggetto che svolazza nell'aria
avrà coordinate come $(-3,-2,5)$. Queste sono \textbf{coordinate globali}
(\emph{world coordinates}), perché sono basate sul \textbf{punto di origine del
mondo} di gioco.

\subsection{Coordinate locali: la spada e la cassa}

Una \textbf{coordinata locale} è una coordinata \textbf{in base al genitore
dell'oggetto}. L'esempio della lavagna: il giocatore ha in mano una
\textbf{spada} (figura~\ref{gd11:fig:coordinate}).
\begin{itemize}
  \item La coordinata \textbf{globale} della spada è la distanza tra
        l'origine del mondo e la spada.
  \item La coordinata \textbf{locale} della spada è la distanza tra la spada
        e l'\textbf{origine del giocatore}, che di solito si trova
        \textbf{ai piedi}, a rappresentare lo zero.
\end{itemize}

\begin{figure}[htbp]
\centering
\begin{tikzpicture}[
  font=\small,
  gd11ax/.style={-{Stealth[length=2.5mm]}, thick},
  gd11w/.style={-{Stealth[length=2.5mm]}, very thick, red!70!black},
  gd11l/.style={-{Stealth[length=2.5mm]}, very thick, blue!70!black, dashed}
]
  % assi mondo
  \draw[gd11ax] (0,0) -- (8.2,0) node[right] {$x$};
  \draw[gd11ax] (0,0) -- (0,4.2) node[above] {$z$ (alto, in Unreal)};
  \fill (0,0) circle (2.2pt);
  \node[below left, font=\footnotesize] at (0,0) {$(0,0,0)_W$};
  % giocatore
  \draw[thick] (5,0) -- (5,1.6);           % corpo
  \draw[thick] (5,1.85) circle (0.25);      % testa
  \draw[thick] (5,0) -- (4.7,-0.0);
  \draw[thick] (5,1.2) -- (5.5,1.0);        % braccio
  \draw[very thick, gray!60!black] (5.5,1.0) -- (6.4,1.6); % spada
  \fill[blue!70!black] (5,0) circle (2.2pt);
  \node[below, font=\footnotesize, blue!70!black] at (5,0) {origine locale (piedi)};
  % frecce
  \draw[gd11w] (0,0) to[bend left=22] (6.0,1.35);
  \node[font=\footnotesize, red!70!black, anchor=south] at (2.6,1.75)
       {coordinata \textbf{globale} della spada};
  \draw[gd11l] (5,0) -- node[right, font=\footnotesize, pos=0.6]
       {\textbf{locale}} (6.0,1.35);
  % cassa
  \draw[thick, fill=brown!20] (1.0,2.6) rectangle (2.6,3.6);
  \fill[orange] (1.4,2.85) circle (0.12);
  \fill[green!50!black] (2.1,3.1) circle (0.12);
  \node[font=\footnotesize, align=left, anchor=west] at (2.8,3.1)
       {la cassa ruota: gli oggetti figli ruotano con lei,\\ma le loro coordinate \textbf{locali} non cambiano};
\end{tikzpicture}
\caption{Coordinate globali e locali, come alla lavagna. La posizione globale
si misura dall'origine del mondo; quella locale dall'origine del genitore (qui il
giocatore, con l'origine ai piedi). Se il genitore si sposta o ruota, cambia la
posizione globale dei figli, non quella locale.}
\label{gd11:fig:coordinate}
\end{figure}

\paragraph{Perché è importante.} Quando si lavora in un livello, ma anche in
qualsiasi cosa di programmazione, bisogna \textbf{rendersi conto se le
coordinate che si stanno usando sono globali o locali}. Il motivo è che gli
oggetti possono essere assegnati come \textbf{figli} di un altro oggetto. La
spada è una cosa diversa dal giocatore, ma si sposta con il giocatore perché
segue le \textbf{trasformazioni} applicate al giocatore. Una
\textbf{trasformazione} (\emph{transform}) è uno \textbf{spostamento}, un
\textbf{cambio di dimensioni} o una \textbf{rotazione}. Se il giocatore si
sposta o si gira, la spada si sposta o si gira; ma \textbf{le sue coordinate
locali non cambiano}: all'interno del sistema di coordinate del giocatore, la
spada è sempre nello stesso punto.

\paragraph{L'esempio della cassa.} Quello della spada ``non è il massimo come
esempio'', perché la spada si muove anche per conto suo, attaccata alla mano. Più
pulito: una \textbf{cassa} con dentro degli oggetti. Se si gira la cassa, ci si
aspetta che si girino anche gli oggetti. Se sono \textbf{figli} della cassa,
si girano con lei; ma all'interno delle coordinate locali gli oggetti \textbf{non
si sono spostati}, ``perché il genitore, il mondo locale, si è spostato''. Lo
zero locale si muove con il genitore; rispetto a quello zero, i figli restano
sempre uguali.

\subsection{La pausa}

Il progetto nuovo deve compilare gli shader da zero (``la speranza è l'ultima a
morire''), quindi il docente anticipa leggermente la \textbf{pausa}: una ventina
di minuti, ``potete chiacchierare, ve lo concedo''. Durante la pausa risponde ad
alcune domande individuali:
\begin{itemize}
  \item \emph{Si possono unire più oggetti in un unico oggetto?} Risposta
        breve: non si può. Risposta lunga: ci saranno lezioni dedicate alla
        \textbf{modellazione 3D}, perché fondere più cose in uno stesso oggetto
        si può fare, ma è più complesso. Per ora: \textbf{raggrupparli}
        (Ctrl+G, sezione~\ref{gd11:sec:gerarchie}).
  \item Gli asset dell'engine non si modificano (sezione~\ref{gd11:sec:asset}).
  \item Per far seguire un oggetto a un altro, nello spazio locale, lo si può
        \textbf{parentare}: lo si trascina sull'altro nell'Outliner, e ``vedete
        che adesso sono uno sotto l'altro''; se lo si sposta, si sposta anche
        l'altro. Rifacendo il trascinamento, li si separa. ``Adesso lo faccio
        vedere a tutti.''
\end{itemize}
Prima di riprendere, l'invito a tenere Unreal aperto, con un progetto
qualsiasi (anche quello del task), per vedere i concetti direttamente
nell'interfaccia.

% ---------------------------------------------------------------------
\section{Lo spazio delle coordinate nell'editor}
\label{gd11:sec:spazio}

Ripresa la lezione (Mattia, da casa, sta finendo di installare la 5.8 sul
portatile), il docente mostra in pratica i concetti spiegati alla lavagna.

\paragraph{Non serve un genitore.} Prima una precisazione: la distinzione tra
coordinate globali e locali vista alla lavagna serviva a dare un'idea di come si
ragiona. In realtà, \textbf{non è necessario che ci sia un genitore} per parlare
di coordinate locali: \textbf{ogni oggetto ha le proprie}.

\paragraph{Il pulsante con il mondo.} Cliccando su un qualsiasi oggetto, in
alto nella viewport c'è un pulsantino con l'icona di un \textbf{mondo}: indica
\textbf{in quale spazio di coordinate} si sta lavorando. Nella traduzione
italiana si chiama \textbf{Spazio del mondo} (``non chiedetemi come sono state
fatte le scelte di traduzione di Epic: sospetto, vista la loro passione per
l'intelligenza artificiale, che c'entri qualcosa''). Cliccandoci sopra si passa
in rassegna:
\begin{description}
  \item[Spazio del mondo] (\emph{World}): le frecce del gizmo indicano gli assi
        $x, y, z$ \textbf{dell'intero mondo di gioco}, ``il punto di
        riferimento più alto che si può avere''. Ruotando un oggetto, le frecce
        restano uguali.
  \item[Spazio locale] (\emph{Local}): le frecce \textbf{seguono gli assi
        dell'oggetto stesso}, e ruotano con lui.
  \item[Spazio genitore] (\emph{Parent}): le frecce seguono gli assi del
        \textbf{genitore} dell'oggetto. Se l'oggetto non è figlio di niente, il
        suo genitore è il mondo, e lo spazio genitore coincide con quello
        globale.
\end{description}
``Ricordatevi: globali, locali, parentali.''

\paragraph{La trave in diagonale.} L'applicazione è chiara. Si prenda un
oggetto allungato e ruotato in diagonale (una \textbf{trave}, un ponte). Se lo
si deve spostare lungo la sua direzione lavorando in coordinate globali, è
``un po' antipatico essere precisi'': si può trascinare il centro del gizmo, ma
ci si sposta leggermente a destra e a sinistra, senza garanzia di precisione.
Passando alle \textbf{coordinate locali}, invece, la trave la si trascina
\textbf{perfettamente nel suo ``avanti''}, lungo il suo asse $X$ locale
(figura~\ref{gd11:fig:trave}). ``È parecchio importante.''

\begin{figure}[htbp]
\centering
\begin{tikzpicture}[
  font=\small,
  gd11ax/.style={-{Stealth[length=2.5mm]}, line width=1.4pt},
  gd11beam/.style={draw, thick, fill=brown!25}
]
  % mondo
  \node[font=\small\bfseries] at (1.6,2.8) {Spazio del mondo};
  \begin{scope}[rotate around={30:(1.6,1.0)}]
    \draw[gd11beam] (0.0,0.8) rectangle (3.2,1.2);
  \end{scope}
  \draw[gd11ax, red!80!black] (1.6,1.0) -- ++(2.1,0) node[right] {$X$};
  \draw[gd11ax, green!50!black] (1.6,1.0) -- ++(0,1.3) node[above] {$Y$};
  \node[font=\footnotesize, align=center] at (1.6,-0.6)
    {le frecce restano allineate al mondo:\\spostare lungo la trave è impreciso};
  % locale
  \node[font=\small\bfseries] at (7.0,2.8) {Spazio locale};
  \begin{scope}[rotate around={30:(7.0,1.0)}]
    \draw[gd11beam] (5.4,0.8) rectangle (8.6,1.2);
    \draw[gd11ax, red!80!black] (7.0,1.0) -- ++(2.1,0) node[right] {$X$};
    \draw[gd11ax, green!50!black] (7.0,1.0) -- ++(0,1.3) node[above] {$Y$};
  \end{scope}
  \node[font=\footnotesize, align=center] at (7.0,-0.6)
    {le frecce ruotano con l'oggetto:\\la trave scorre nel suo ``avanti''};
\end{tikzpicture}
\caption{Lo stesso oggetto ruotato, con il gizmo in spazio del mondo e in spazio
locale (vista dall'alto). La \emph{scala}, invece, lavora sempre in spazio
locale.}
\label{gd11:fig:trave}
\end{figure}

% ---------------------------------------------------------------------
\section{Genitori, figli e gruppi}
\label{gd11:sec:gerarchie}

\subsection{Parent e child}

Lo spazio genitore introduce il concetto di \textbf{genitorialità}: il rapporto
\textbf{parent} e \textbf{child}, genitore e figlio. Sullo schermo, due cubi:
quello a sinistra è il genitore, quello a destra il figlio.

\paragraph{Come si crea.} Nell'\textbf{Outliner}, a destra
(sezione~\ref{gd8:sec:outliner}), si afferra un oggetto e lo si \textbf{trascina
su un altro}: lo si sta assegnando come suo figlio. I figli compaiono sotto il
genitore, e si possono \textbf{racchiudere e riaprire} con la freccina accanto
al nome (il docente rinomina il figlio \emph{Cube Child}).

\paragraph{L'asterisco.} Una piccola parentesi: l'\textbf{asterisco} accanto al
nome di un livello o di un asset significa che \textbf{c'è qualcosa che non è
stato salvato}. Spostando un cubo, l'asterisco compare: il livello è stato
modificato. ``Se chiudete senza salvare, tutte le robe con gli asterischi sono
perse.'' \textbf{Ctrl+S} salva; \textbf{Ctrl+Shift+S} salva tutto
(sezione~\ref{gd8:sec:livelli}).

\paragraph{Che cosa eredita un figlio.} ``Che cos'è un child? Un child è
qualcosa che \textbf{eredita tutte le trasformazioni del genitore}.'' Il docente
chiede a una volontaria (``mi sembra una faccia da volontaria'') che cosa
succede se si sposta il genitore: ``si sposta anche l'altro''. Esatto: il figlio
eredita, ``da qui il termine'', tutte le trasformazioni applicate al genitore,
\textbf{incluse quelle di grandezza}. E qui c'è un dettaglio da notare: quando
si \textbf{scala} il genitore, non si allarga solo l'oggetto, ma anche la
\textbf{distanza del figlio dallo zero locale}, cioè dall'origine del genitore.
Se si stretcha il genitore, si stretcha \textbf{l'intero insieme}, non i singoli
oggetti ciascuno per conto proprio. Lo stesso per la rotazione e lo spostamento
(e Ctrl+Z, Ctrl+Z, Ctrl+Z per tornare indietro).

\paragraph{Le gerarchie a più livelli.} Con Ctrl+C e Ctrl+V (o Ctrl+D) si
duplica il figlio, e si assegna il nuovo cubo (\emph{Cube Child 2}) come figlio
di \emph{Cube Child}: ora c'è una \textbf{discendenza}
(figura~\ref{gd11:fig:gerarchia}). Se le cose sono ereditate a livello di
gerarchia, \textbf{quello più in alto influenza tutti quelli che gli stanno
sotto}; ma modificando il secondo si influenza solo il suo figlio.

\begin{figure}[htbp]
\centering
\begin{tikzpicture}[
  font=\small,
  gd11node/.style={draw, thick, rounded corners=3pt, fill=blue!8,
                   minimum width=2.6cm, align=center, inner sep=4pt},
  gd11grp/.style={draw, thick, rounded corners=3pt, fill=orange!12,
                  minimum width=2.6cm, align=center, inner sep=4pt},
  gd11arr/.style={-{Stealth[length=2.3mm]}, thick}
]
  \node[font=\small\bfseries] at (1.5,0.9) {Gerarchia (parent/child)};
  \node[gd11node] (c0) at (1.5,0) {\emph{Cube}};
  \node[gd11node] (c1) at (2.3,-1.2) {\emph{Cube Child}};
  \node[gd11node] (c2) at (3.1,-2.4) {\emph{Cube Child 2}};
  \draw[gd11arr] (c0.south west) ++(0.3,0) |- (c1.west);
  \draw[gd11arr] (c1.south west) ++(0.3,0) |- (c2.west);
  \node[font=\footnotesize, align=left, anchor=north west] at (0.2,-3.0)
    {muovo \emph{Cube}: si muovono tutti\\muovo \emph{Cube Child}: si muove anche \emph{Child 2}\\muovo \emph{Child 2}: si muove solo lui};
  \node[font=\small\bfseries] at (8.6,0.9) {Gruppo (Ctrl+G)};
  \node[gd11grp] (g) at (8.6,0) {\emph{GroupActor}};
  \node[gd11node, minimum width=1.4cm] (a) at (7.2,-1.4) {A};
  \node[gd11node, minimum width=1.4cm] (b) at (8.6,-1.4) {B};
  \node[gd11node, minimum width=1.4cm] (c) at (10.0,-1.4) {C};
  \foreach \n in {a,b,c} \draw[thick, dashed] (g) -- (\n);
  \node[font=\footnotesize, align=left, anchor=north west] at (6.6,-2.2)
    {asset separati, trattati come un'unità:\\cliccandone uno si selezionano tutti;\\ideale per spostare e ruotare\\(es. i libri di una libreria)};
\end{tikzpicture}
\caption{Gerarchia e gruppo. Nella gerarchia ogni livello eredita le
trasformazioni di quello sopra; nel gruppo gli oggetti restano indipendenti ma
vengono selezionati e trasformati insieme.}
\label{gd11:fig:gerarchia}
\end{figure}

\subsection{I gruppi}

``La differenza, che potrebbe essere una domanda, tra questo e un
\textbf{gruppo}.'' Selezionando più oggetti e premendo \textbf{Ctrl+G}, Unreal
li raggruppa e crea un \textbf{Group Actor}, che segnala che tutti quegli
oggetti formano un \textbf{gruppo unico}. Sono asset separati, ma l'editor dice:
``queste cose qua sono un unico gruppo, considerale insieme''.
\begin{itemize}
  \item Spostando il gruppo, \textbf{si spostano tutti insieme}.
  \item Selezionando \textbf{uno} degli oggetti, il Group Actor li fa
        selezionare \textbf{tutti}, anche nell'Outliner.
  \item Si possono anche \textbf{scalare} insieme, come se fossero figli, anche
        se non lo sono; ma ``se stretchate una roba raggruppata, possono
        succedere cose strane''. Idealmente, i gruppi servono a
        \textbf{spostare e ruotare} cose insieme.
\end{itemize}
Il caso classico che capiterà di vedere nelle mappe fatte da altri è la
\textbf{libreria}: i libri sono raggruppati, perché così è comodo spostarli e
piazzarli in giro.

\paragraph{Stretchare: quasi solo nel blockout.} Una nota a margine: lo
\textbf{stretch} (la scala non uniforme) è qualcosa che capiterà quasi solo in
fase di \textbf{blockout}. Una spada stretchata si vede che è stretchata, un
mobile stretchato si vede che è stretchato: ``è difficile che vi capiterà di
stretchare le cose nella versione finale di un ambiente'' (si ricordi il
problema dello stretch nella sezione~\ref{gd11:sec:spline}).

\paragraph{Una domanda dal fondo dell'aula.} ``Ci siete? Sto giocando a
\emph{Geometry Dash} là dietro: che state facendo?'' Va bene, si riprende.

% ---------------------------------------------------------------------
\section{Ctrl+Z: anche cliccare è un'azione}
\label{gd11:sec:ctrlz}

``Piccola nota importante, che probabilmente avete già notato, ma è meglio
dirlo in maniera esplicita.'' \textbf{Ctrl+Z} annulla l'ultima azione. Ma in
Unreal, e in molti altri editor 3D, \textbf{cliccare su qualcosa nella viewport
è considerato un'azione}. Se si clicca un oggetto, poi un altro, poi un altro, e
si preme Ctrl+Z, si torna alla \textbf{selezione precedente}: non si annulla
l'ultima modifica.

È utile, perché permette di tornare indietro anche nelle selezioni; ma
all'inizio è \textbf{controintuitivo}, perché ci si aspetta che Ctrl+Z annulli
l'ultima cosa \emph{fatta}, non l'ultima cosa \emph{selezionata}. Il rischio
concreto: si fanno tre modifiche, si pensa ``no, meglio di no'', si preme Ctrl+Z
tre volte; ma anche i clic tra una modifica e l'altra contano come azioni da
annullare, e le modifiche non sono state annullate tutte. ``Ricordate questa
cosa: \textbf{cliccare su un oggetto è un'azione}.''

% ---------------------------------------------------------------------
\section{Lo snapping}
\label{gd11:sec:snapping}

\subsection{Che cos'è}

``Che cos'è lo snapping? È molto semplice.'' Se si ha una \textbf{griglia} e si
stanno posizionando degli oggetti (per esempio una serie di piattaforme), la
griglia fa sì che ogni oggetto \textbf{si incastri automaticamente} in una
posizione precisa. Dall'inglese \emph{to snap}: qualcosa che \textbf{scatta} in
una posizione. Afferrando un oggetto e trascinandolo, l'editor lo piazza
automaticamente nel punto giusto della griglia. È molto utile con i
\textbf{kit modulari} (sezione~\ref{gd11:sec:modulari}), ma ``per un sacco di
altra roba''.

\subsection{La barra degli strumenti della viewport}

In alto a destra nella viewport ci sono i controlli già visti
(sezione~\ref{gd8:sec:transform}): la \textbf{selezione} (che ``non usa mai
nessuno da sola'', tanto vale trascinare direttamente), lo spostamento, la
\textbf{rotazione}, la \textbf{scala} (W, E, R), e lo \textbf{spazio delle
coordinate} (sezione~\ref{gd11:sec:spazio}). Precisazione: la \textbf{scala è
sempre locale}.

\paragraph{La griglia di posizione.} Poi c'è il simbolino della
\textbf{griglia}, ``presente in tutti gli editor 3D''. Spostando un cubo, si nota
che non si muove di un pixel alla volta, ma a \textbf{incrementi di 10}: 10
unità di Unreal, cioè \textbf{10 centimetri} (``le unità di Unreal sono uno a
uno con i centimetri''), perché la griglia è impostata a 10.

\paragraph{Il cubo da un metro.} Il \textbf{cubo di base} di Unreal è largo
\textbf{un metro}: un metro per un metro per un metro, con il pivot al centro
(sezione~\ref{gd11:sec:pivot}). ``Un metro, quanti centimetri sono?'' (``Due
tentativi, ma ce l'abbiamo fatta'': cento.) Impostando la griglia a
\textbf{100}, il cubo si sposta esattamente della propria larghezza. Con
Ctrl+C e Ctrl+V si possono affiancare copie perfettamente allineate e costruire
figure molto velocemente.

\paragraph{I valori disponibili.} I valori della griglia sono \textbf{limitati}
a quelli proposti nel menu (``un po' antipatico''): da qualche parte c'è un modo
per impostarli a piacere, ma se si mantengono buone dimensioni quelli presenti
bastano.

\paragraph{Quando spegnerla.} Se invece si stanno mettendo, per esempio, delle
\textbf{rocce} o degli \textbf{oggetti su un tavolo}, e gli scatti della griglia
danno un risultato ``molto falso'', poco naturale, si clicca sull'icona per
\textbf{disattivarla} e si sposta l'oggetto liberamente. Il consiglio:
\begin{quote}
Usate la griglia quando serve. Abituatevi a pensare: questa cosa è più comoda se
la faccio con una griglia che scatta, o no? Vi velocizza tanto il lavoro,
rispetto a stare lì manualmente a spostare le cose un centimetro alla volta.
\end{quote}

\subsection{Lo snapping della rotazione}

Lo stesso vale per la \textbf{rotazione}: di default procede a scatti di
\textbf{10 gradi}. Disattivando lo snap si ruota liberamente. Per ruotare
qualcosa esattamente di \textbf{90 gradi} si imposta 90, ``tac'' (``è un cubo,
non si vede la differenza, però fidatevi''); per 45 gradi, 45.

\paragraph{I valori a mano, e le operazioni.} Per un valore preciso che non è
tra quelli proposti, lo si scrive nel campo della rotazione nei
\emph{Details}: ruotare di 12 gradi sull'asse $z$? Si scrive 12. Ruotarlo tre
volte tanto? Si scrive \texttt{12*3}, Invio, e il risultato è 36 (il trucco dei
campi numerici, sezione~\ref{gd8:sec:transform}). ``La matematica, purtroppo,
sarà estremamente importante in questo tipo di mestiere.''

\subsection{Ruotare attorno a un asse}

``Qualcuno si è chiesto: perché l'ha ruotato sull'asse $z$?'' L'asse $z$ è
quello \textbf{verticale}. Quando si ruota \textbf{su un asse}, bisogna
immaginare che si ruoti \textbf{attorno} a quell'asse: l'asse fa da
\textbf{perno}. La rotazione sull'asse $z$ (che infatti è blu, come la freccia
di $z$) fa girare l'oggetto \textbf{orizzontalmente}, come una trottola.

Anche questo all'inizio può essere controintuitivo: si pensa che per farlo
ruotare orizzontalmente si debba usare l'asse $x$, ma ruotando sull'asse $x$ lo
si fa girare \textbf{attorno all'asse orizzontale}, non sul piano orizzontale
(tabella~\ref{gd11:tab:assi}).

\begin{table}[htbp]
\centering
\small
\begin{tabular}{@{}llll@{}}
\toprule
\textbf{Asse (Unreal)} & \textbf{Colore} & \textbf{Direzione} & \textbf{Ruotando attorno ad esso\ldots} \\
\midrule
$X$ & rosso & avanti/indietro & l'oggetto si inclina di lato (\emph{roll}) \\
$Y$ & verde & destra/sinistra & l'oggetto si inclina avanti/indietro (\emph{pitch}) \\
$Z$ & blu & alto/basso & l'oggetto gira sul piano orizzontale (\emph{yaw}) \\
\bottomrule
\end{tabular}
\caption{Gli assi di Unreal e l'effetto della rotazione attorno a ciascuno. I
nomi \emph{roll}, \emph{pitch} e \emph{yaw} sono quelli usati nei campi di
rotazione dell'editor; a lezione il docente insiste soprattutto sull'asse $z$
come perno verticale.}
\label{gd11:tab:assi}
\end{table}

% ---------------------------------------------------------------------
\section{Il Content Browser e la lingua dell'editor}
\label{gd11:sec:browser}

\paragraph{Il cassetto dei contenuti.} ``Piccola altra cosa'': il
\textbf{Content Browser}, che nella versione italiana si chiama
\textbf{Cassetto dei contenuti}. Il docente usa Unreal in inglese e si scusa se
non usa sempre la terminologia italiana. Il Content Browser si può
\textbf{chiudere}: premendo \textbf{Ctrl+Spazio} lo si fa apparire, premendolo
di nuovo lo si nasconde. Così si ha più spazio per vedere il mondo di gioco.

\paragraph{Cambiare la lingua.} Uno studente chiede come cambiare la lingua
dell'editor. Il docente non è sicuro che il percorso non sia cambiato nella 5.8,
ma dovrebbe essere: \textbf{Modifica} $\rightarrow$ \textbf{Impostazioni
dell'editor} (\emph{Edit} $\rightarrow$ \emph{Editor Preferences}), e cercare
``lingua'': lì la si sceglie. Da non confondere con l'\textbf{anteprima della
lingua del gioco}, che serve, nei giochi localizzati in più lingue, a scegliere
in quale lingua giocarli quando li si fa partire. Il docente, come si è visto,
preferisce l'inglese: è più comodo, tra l'altro, per fare \textbf{ricerche}
(sezione~\ref{gd11:sec:viewmode}).

% ---------------------------------------------------------------------
\section{Il pivot}
\label{gd11:sec:pivot}

\slidefig{img/GD/gd11-pivot}{Pivot, mesh e luci.}{slide \emph{Intro Level Design Unreal} (2)}{fig:gd11-pivot}


Si torna alle slide (``non chiudete Unreal, ci servirà''). \textbf{Pivot} (``a
seconda che scegliate di pronunciarlo alla francese o all'inglese''): il pivot è
\textbf{il punto che un actor considera l'origine delle sue coordinate}. Il cubo
di base di Unreal ha il pivot \textbf{esattamente al centro}. Quindi, se lo si
ruota, ruota \textbf{attorno al suo centro}. (Si era già visto, nella lezione
precedente, che non tutti gli oggetti hanno il pivot al centro: lo decide chi
crea la mesh, sezione~\ref{gd8:sec:editor} e figura~\ref{gd8:fig:gizmo}.)

\paragraph{Il tasto G.} Parentesi utile, anch'essa nella slide con le
scorciatoie: premendo \textbf{G} si \textbf{nascondono i gizmo} e tutti gli
strumenti interni dell'editor (frecce, righe, icone), per vedere la scena com'è
davvero. Premendolo di nuovo, tornano.

\paragraph{Ruotare attorno a un angolo.} ``Molto figo, ma se questo cubo lo
volessi ruotare attorno a un angolo?'' Il caso classico: si sta posizionando un
\textbf{ponte}, si sa dov'è il suo punto di partenza, ma il pivot del ponte è al
centro, e ruotarlo per farlo combaciare è ``una rottura'': ruotare, spostare,
non va bene, ruotare, spostare. Sarebbe comodo prendere il pivot e spostarlo
all'inizio del ponte, così da ruotarlo attorno al punto di partenza.

\paragraph{Alt + tasto centrale.} Si può fare: tenendo premuto \textbf{Alt} e
cliccando con il \textbf{tasto centrale del mouse}, si sposta il pivot
dell'oggetto \textbf{nel punto in cui si clicca}. Il docente clicca su un angolo
del cubo, e ora il cubo ruota attorno a quell'angolo. Il punto può anche essere
\textbf{esterno} all'oggetto: con Alt e clic centrale fuori dal cubo, il pivot
finisce nel vuoto, e il cubo ruota (o si trascina) a partire da quel punto
(figura~\ref{gd11:fig:pivot}). Può tornare utile per certi tipi di
spostamenti.

\begin{figure}[htbp]
\centering
\begin{tikzpicture}[
  font=\small,
  gd11cube/.style={draw, thick, fill=gray!20},
  gd11ghost/.style={draw, thick, dashed, gray},
  gd11piv/.style={circle, fill=red!75!black, inner sep=2pt},
  gd11rot/.style={-{Stealth[length=2.2mm]}, thick, blue!70!black}
]
  \foreach \cx/\title/\px/\py in {1.2/{pivot al centro}/1.2/1.0,
                                   5.4/{pivot su un angolo}/4.9/0.5,
                                   9.8/{pivot esterno}/10.9/1.0} {
    \node[font=\small\bfseries] at (\cx,2.6) {\title};
    \draw[gd11ghost] (\cx-0.5,0.5) rectangle (\cx+0.5,1.5);
    \begin{scope}[rotate around={35:(\px,\py)}]
      \draw[gd11cube] (\cx-0.5,0.5) rectangle (\cx+0.5,1.5);
    \end{scope}
    \node[gd11piv] at (\px,\py) {};
    \draw[gd11rot] (\px,\py) ++(0:0.55) arc (0:120:0.55);
  }
  \node[font=\footnotesize, align=center] at (5.4,-1.3)
    {Alt + clic centrale sposta il pivot \textbf{temporaneamente}\\(anche fuori dall'oggetto); deselezionando torna quello del modellatore};
\end{tikzpicture}
\caption{Lo stesso cubo ruotato di 35 gradi attorno a tre pivot diversi (in
rosso). Il tratteggio indica la posizione di partenza.}
\label{gd11:fig:pivot}
\end{figure}

\paragraph{Uno spostamento temporaneo.} Questo spostamento è
\textbf{temporaneo}. Se non ci si ricorda la scorciatoia, lo si può fare anche
dal \textbf{tasto destro} del mouse sull'oggetto, dove c'è una voce dedicata al
pivot (per esempio \emph{Imposta come offset del pivot}: in aula non funziona al
primo colpo, ma le opzioni ci sono). Ci sono modi per \textbf{bloccare} il pivot
e per impostare un pivot alternativo \textbf{che non viene resettato}; ma per
adesso è meglio tenere a mente che \textbf{si resetta ogni volta}, ``così non
rischiate di fare casini'': si sa sempre che il pivot è nel punto in cui l'ha
impostato il \textbf{modellatore} che ha creato l'oggetto, e lo si cambia
temporaneamente con Alt + tasto centrale solo quando serve.

% ---------------------------------------------------------------------
\section{Mesh e collisioni}
\label{gd11:sec:mesh}

\subsection{Che cos'è una mesh}

``Una nota importante: le mesh.'' Detto ``terra terra, nella maniera più
brutale possibile'', una \textbf{mesh} è \textbf{un insieme di punti e di
segmenti} che connettono questi punti. (Il docente la disegna alla lavagna:
``i modellatori di Tadic che vedono che disegno una roba del genere mi
sparano''\footnote{Il nome è riportato così nella trascrizione; non è chiaro a
chi o a quale studio si riferisca.}, ``ma è illustrativo''.)

\paragraph{Tutto è fatto di triangoli.} Chiunque abbia giocato a un videogioco
avrà notato che, guardando le cose da vicino, sono composte di
\textbf{poligoni}. ``Non stiamo a parlare di quad e tutto il resto: stiamo
semplici sulla teoria della modellazione 3D. Si parla di \textbf{triangoli}.''
Perché il triangolo? Perché è \textbf{la figura geometrica che si può costruire
con il minor numero di punti}: un punto non è una figura, una linea non ha
superficie. Tre punti collegati formano la \textbf{superficie più semplice}
possibile: il triangolo è \textbf{l'unità di base} della grafica 3D. ``In
maniera molto semplificata, perché non è esattamente vero, ma per i nostri
scopi'': tutto ciò che si vede in un engine 3D è fatto di piccolissimi
triangoli.

\paragraph{Una rete.} \emph{Mesh}, in inglese, vuol dire \textbf{rete},
``retina'', come nei vestiti a rete. Una mesh è un \textbf{file} che specifica
\textbf{l'insieme dei punti} (vertici) e \textbf{in che modo sono collegati tra
loro} (le linee); le superfici delimitate da queste linee sono i
\textbf{poligoni}, su cui vengono poi proiettati i \textbf{materiali} e le
\textbf{texture} (sezione~\ref{gd11:sec:materiali}).

\paragraph{La testa in wireframe.} Sulla slide c'è la mesh di un corpo: la
testa è formata da poligoni. Quelli visualizzati non sono tutti triangoli (molti
sono quadrilateri), ma ``non stiamo entrando nel discorso di come funziona la
modellazione 3D'': edge, raggruppamenti e così via. Lo si potrà vedere in
Unreal con la modalità \emph{wireframe} (sezione~\ref{gd11:sec:viewmode}).

\subsection{Le collisioni sono un'altra cosa}

``Nota importantissima, nota estremamente importante'':
\begin{quote}
\textbf{Mesh e collisioni sono due cose diverse.} Se vedete qualcosa in 3D, non
è detto che sia fisico: qualcos'altro deve dire all'engine che quella cosa è
solida.
\end{quote}
Le mesh sono solo una \textbf{rappresentazione tridimensionale virtuale}. La
\textbf{collisione} è una cosa completamente separata, e \textbf{non è
automatica}. Unreal e molti altri engine possono usare la mesh stessa come
modello di collisione, ma non è automatico e \textbf{non è consigliato}, perché
la mesh è super dettagliata e il calcolo diventa molto pesante (di solito si usa
una forma di collisione più semplice).

\paragraph{Lo studente che sprofonda.} Quello che deve restare in testa:
mettendo in gioco un modello 3D, \textbf{non è automatico che abbia una
collisione}. Bisogna \textbf{controllarlo}, perché ``è una cosa che succede
spesso agli studenti'': ``che bello questo livello''; il personaggio cammina,
salta sulla roccia e \textbf{sprofonda nel terreno}, perché la roccia non ha la
collisione.

% ---------------------------------------------------------------------
\section{Le luci}
\label{gd11:sec:luci}

``Altra cosa importante da vedere sono le luci. Per questo mi farebbe comodo un
po' di buio.'' Ogni nuovo livello creato da un template di Unreal ha una
cartella \textbf{Lighting}, che si vedrà tra poco
(sezione~\ref{gd11:sec:da-zero}); per ora il docente abbassa l'intensità della
\emph{directional light} per avere una scena più scura. Nel pannello a sinistra
(\emph{Place Actors}) c'è la categoria \textbf{Luci}: sono tutte le luci che si
possono inserire (tabella~\ref{gd11:tab:luci}). Il dettaglio di ciascuna verrà
in lezioni dedicate: ``questo è per farvi vedere cosa sono'' e mettervi in mano
la possibilità di fare cose velocemente.

\begin{table}[htbp]
\centering
\small
\begin{tabular}{@{}p{2.6cm}p{4.2cm}p{5.6cm}@{}}
\toprule
\textbf{Luce} & \textbf{Come emette} & \textbf{Note dalla lezione} \\
\midrule
Point light (punto luce) & in ogni direzione, a sfera &
  raggio di attenuazione; \emph{source radius}: se maggiore di zero la luce non
  è più un punto e le ombre si sfumano; colore, intensità \\
\addlinespace
Spot light (faretto) & a cono, in una direzione &
  come un faro, una torcia; intensità in candele (cd), il cursore si ferma a
  un massimo ma a mano si può andare oltre \\
\addlinespace
Rect light (luce rettangolare) & da una superficie, a semisfera &
  poco usata; utile per una finestra illuminata che proietta luce verso
  l'esterno \\
\addlinespace
Directional light (luce direzionale) & raggi paralleli e infiniti su tutto il mondo &
  è il sole: conta solo la rotazione, non la posizione; legata al cielo
  (tramonto); \textbf{una sola per livello} \\
\bottomrule
\end{tabular}
\caption{Le luci di Unreal viste a lezione.}
\label{gd11:tab:luci}
\end{table}

\subsection{La point light}

Si comincia dal \textbf{punto luce} (\emph{point light}). Trascinandolo nel
livello (il docente sposta alcuni oggetti sul terreno, per avere qualcosa su cui
proiettare), si vede che \textbf{all'interno di una sfera}, il gizmo della luce,
viene proiettata luce.

\paragraph{La luce non è un'entità.} Per selezionarla si clicca
sull'\textbf{iconcina}: una luce non è un oggetto, ``è un punto di diramazione di
raggi di luce''. Le iconcine servono a due cose: \textbf{segnalare} che lì c'è
qualcosa, e permettere di \textbf{cliccarci sopra}. Premendo G (sezione
precedente), infatti, lì non si vede niente.

\paragraph{Il modello della lampadina non fa luce.} Ed ecco un'altra cosa
importante: \textbf{la fonte di luce visibile} (il modello 3D di una lampadina,
di un fuoco) \textbf{non emette luce} se non le si mette accanto una luce.
``Poi in realtà un pochino ne emette, perché ci sono riflessi di colori
ambientali, ma quella è un'altra cosa.'' Le luci dell'engine sono solo la
\emph{sorgente}; il modello è solo l'aspetto.

\paragraph{Il raggio.} La luce si può selezionare e spostare, e spostandola si
spostano anche le \textbf{ombre} che proietta. Ma le luci \textbf{non hanno
influenza all'infinito} (``altrimenti buona fortuna a far funzionare tutto''):
sono limitate al \textbf{raggio} visualizzato. Gli oggetti fuori dal raggio non
vengono toccati dalla luce.

\paragraph{Le proprietà.} Nei \emph{Details} ci sono tutte le funzioni per
modificarla:
\begin{itemize}
  \item il \textbf{colore}, con il selettore;
  \item il \textbf{raggio di attenuazione} (\emph{attenuation radius}): la
        distanza oltre la quale l'influenza della luce su un oggetto diventa
        zero;
  \item la \textbf{dimensione della sorgente} (\emph{source radius}): se è
        zero, la sorgente è un \textbf{punto}; se è più larga di zero, è la
        dimensione della sfera da cui viene proiettata la luce. Allargandola, le
        \textbf{ombre diventano più sfumate}, perché la luce non è più un punto
        ma una sfera.
\end{itemize}

\subsection{La spot light}

La point light proietta luce \textbf{in ogni direzione}, a sfera. La
\textbf{spot light}, invece, è una proiezione \textbf{a cono}: ``lo dice il
termine'', è come un proiettore, ``un faro''. Serve quando invece di mandare la
luce in tutte le direzioni la si vuole lanciare solo in una: ``immaginatela come
una \textbf{torcia}''. Con G si vede il cono di luce proiettato sulla scena;
anche qui la si seleziona cliccando sull'iconcina, e se ne regola tutto,
compresa l'\textbf{intensità}: ``volete spaccare gli occhi alla gente? Mettetela
al massimo''.

\paragraph{I campi non sono sempre bloccati.} Una nota di usabilità.
Trascinando il cursore dell'intensità, si nota che è \textbf{bloccato} a un
massimo (nell'esempio, 160 candele: \textbf{cd}, perché Unreal usa unità di
misura del mondo reale), e tornando indietro si ferma a zero. Ma questi campi
\textbf{non sono sempre limitati brutalmente} ai valori del cursore:
cliccandoci sopra e scrivendo il valore a mano, \textbf{in certi campi si può
andare oltre il limite}. ``Volete una luce da 2000 candele? Sì, però tenete
presente che accecate chiunque.''

\subsection{La rect light}

La \textbf{luce rettangolare} (\emph{rect light}) è una \textbf{superficie} che
proietta luce, non un punto: ruotandola si vede che la luce si interrompe,
perché non proviene da un punto, ma da una superficie che illumina ``a
semisfera''. Non è usata molto spesso, ma ha i suoi scopi: per esempio una
\textbf{finestra illuminata} che proietta luce verso l'esterno di un palazzo,
lungo tutta la sua larghezza.

\subsection{La directional light: il sole}

``Directional light: \textbf{non trascinatela} nel vostro progetto.'' Quando si
apre un livello da un template di Unreal c'è già, nella cartella
\textbf{Lighting} (``ignoriamo la fog per ora, che è la nebbia''). Il docente
la riaccende e la rimette a \textbf{2 lux}.

\paragraph{Che cos'è.} La \textbf{directional light} è \textbf{la luce del
sole}: ``è il vostro sole''. Si chiama \emph{direzionale} perché \textbf{tutti i
raggi di luce sono paralleli}, nella stessa direzione. E se sono tutti paralleli
e infiniti, perché coprono tutto il mondo di gioco, \textbf{il punto di origine
non cambia assolutamente niente}: la si può mettere dove si vuole. Di solito,
una volta trovato un posto in cui non dà fastidio, non la si tocca quasi più.

\paragraph{Il tasto F.} Per ritrovarla, un'altra scorciatoia ``importantissima''
della slide: il tasto \textbf{F} porta la telecamera sull'\textbf{oggetto
selezionato}. Si seleziona qualcosa nell'Outliner, F, e la viewport lo inquadra;
si seleziona un'altra cosa, F, e così via.

\paragraph{Ruotare il sole.} Per cambiare la posizione del sole \textbf{la si
ruota}: ruotandola si spostano anche le ombre. (Il docente passa alle coordinate
locali, ``così ce l'ho più comodo''; e ripete: se un oggetto non è figlio di
niente, il genitore è il mondo, e lo spazio genitore coincide con quello
globale.)

\paragraph{Il sole e il cielo.} Unreal \textbf{lega automaticamente} la
directional light al \textbf{sole} del cielo (si veda la \emph{sky atmosphere},
sezione~\ref{gd11:sec:da-zero}): ruotando la luce si sposta il sole, con tutti
i calcoli che ne derivano, incluso l'aspetto del gioco al \textbf{tramonto},
quando la luce diventa più rossa. ``In questo modo potete giocare a essere Dio e
spostare il sole, le stelle, eccetera.'' Le stelle sono un discorso diverso,
perché quando la luce direzionale è sotto l'orizzonte non c'è niente che
illumina il resto. E una curiosità: \textbf{nei videogiochi la luna, in realtà,
è un sole con una texture diversa}.

\paragraph{Una sola.} ``Non mettetene mai due.'' \textbf{Ogni livello deve avere
una sola directional light}: con due, si costringe l'engine a fare il doppio dei
calcoli, senza senso.

% ---------------------------------------------------------------------
\section{Lo Starter Content, a mano}
\label{gd11:sec:starter}

\subsection{Un nuovo progetto, e un engine cambiato}

``Ora costruiamo un livello da zero.'' Le slide erano state scritte basandosi
sulla versione 5.6, quindi ``ci sono un paio di cose che ignorate pure''. Le
istruzioni erano: creare un \textbf{nuovo progetto} con il \textbf{Third Person
template}. Ma nella 5.8 la schermata di creazione del progetto è cambiata
``totalmente'': il docente crea un progetto vuoto e scopre che non ha più i
preset di base che c'erano una volta (``fantastico''), mentre sui computer degli
studenti i moduli necessari risultano già installati. È solo il computer nuovo a
non averli.

\paragraph{Lo Starter Content non si importa più con un clic.} Il problema
principale è un altro, e il docente l'ha verificato il giorno prima: una volta
bastava fare tasto destro, \emph{importa Starter Content}, e si era finito.
``Hanno tolto sta roba, rottura gigantesca.'' Per questo nella cartella
condivisa c'è un \textbf{link a un pacchetto} con tutto lo Starter Content
(texture, materiali, particelle, prop, eccetera) da importare manualmente.

\paragraph{L'accesso alla cartella.} Gli studenti cominciano a scaricarlo. Il
docente rende il link accessibile a chiunque lo abbia, chiedendo di \textbf{non
richiedere l'accesso}: ogni richiesta di accesso a una cartella manda una mail a
tutti quelli che fanno parte dell'organizzazione, e ``ogni giorno ricevo un
sacco di mail di gente che chiede accesso a roba che non ho mai neanche saputo
che esistesse''. (E il Wi-Fi dell'aula è un po' lento.)

\paragraph{La documentazione.} Intanto, uno sguardo alla pagina da cui si è
scaricato il pacchetto: il link sopra è alla \textbf{documentazione di Unreal}
(nella versione 5.9: a sinistra si può scegliere la versione dell'engine,
``perché un po' di robe cambiano, come avete visto''), che spiega come si
importano le varie cose. Si ricordi l'invito della lezione precedente a leggere
la documentazione (sezione~\ref{gd8:sec:documentazione}).

\subsection{Perché l'importazione normale non funziona}

Il docente prova a importare il pacchetto con \emph{Importa nella cartella
corrente}, ma Unreal non lo prende. Chiede aiuto a uno studente (``Alessandro,
vuoi fare da assistente? Ricordami come si importano gli asset da un progetto
all'altro, che non lo faccio da una vita''), contatta il collega che gli ha
passato il file, e si prende una ``piccola seconda pausa'' per risolvere.

La causa: i file del pacchetto sono \textbf{\texttt{.uasset}}, cioè
\textbf{asset già importati} in Unreal. L'importazione normale serve a
trasformare file esterni (immagini, modelli) in asset; questi asset lo sono già,
e non vanno importati, ma \textbf{messi direttamente nella cartella del
progetto}. (La slide con la procedura di importazione standard era sbagliata,
perché il docente aveva frainteso che tipo di file gli venisse passato: ``ignorate
questa slide''. La documentazione resta utile, ma per importare asset
\emph{non} pre-importati.)

\subsection{La procedura}

La procedura corretta, che gli studenti eseguono in aula:
\begin{enumerate}
  \item \textbf{scaricare} il pacchetto e \textbf{decomprimerlo};
  \item in Unreal, nel Content Browser, \textbf{tasto destro} sulla cartella
        \textbf{Contenuto} (\emph{Content}) e scegliere \textbf{Mostra in
        Esplora risorse} (\emph{Show in Explorer}); in alternativa, aprire a
        mano la cartella \texttt{Content} del progetto;
  \item \textbf{trascinare} la cartella decompressa \textbf{dentro} quella
        cartella, in Windows (non nel Content Browser: ``non così, perché così
        vedete che non importa'');
  \item \textbf{chiudere e riaprire} Unreal.
\end{enumerate}
``Che cosa sta succedendo? State mettendo dei nuovi file di Unreal dentro
Unreal.'' Quando si avvia, Unreal \textbf{scansiona tutti i contenuti del
progetto}, nota che ci sono dei file nuovi e li aggiunge. ``Lo so, era molto più
comodo quando lo si poteva selezionare con un clic.''

\paragraph{Un cuscinetto di minuti.} Il docente si scusa per il caos (``tra il
computer nuovo e il fatto che io non uso la 5.8 per lavoro, incidenti di
percorso'') e spiega che, quando prepara una lezione, si lascia sempre
\textbf{qualche minuto di cuscinetto} per imprevisti come questo, soprattutto
con un engine aggiornato che ha cambiato un bel po' di cose, anche
nell'interfaccia. Di conseguenza può capitare di finire la lezione un po' prima;
in quei casi, se ci sono domande o argomenti di cui parlare, ``chiedetelo voi''.
Mattia, da casa, avvisa che probabilmente riuscirà solo a seguire; ``va
benissimo''.

\subsection{Che cosa c'è dentro}

Intanto si esplora il pacchetto. Lo Starter Content è suddiviso in
\textbf{sottocartelle}:
\begin{itemize}
  \item una cartella che ripete varie \textbf{forme di base} di Unreal (``ci
        interessa poco'');
  \item \textbf{Architecture}: forme di base pensate per fare il blockout di
        architetture varie; ``volete creare velocemente un edificio? Pam pam
        pam, e l'avete fatto'';
  \item \textbf{Props}: oggetti di arredamento, sedie, lampade, lampade da
        parete, e altro;
  \item elementi \textbf{naturali}: una roccia, un cespuglio, cose che fanno
        sembrare tutto ``un po' più naturale'';
  \item una \textbf{statuetta di vetro}, trasparente: il vetro
        \textbf{deforma} ciò che c'è dietro, perché ha un materiale (uno
        \emph{shader}) che deforma lo sfondo;
  \item e soprattutto i \textbf{materiali} (sezione~\ref{gd11:sec:materiali}).
\end{itemize}
\paragraph{Le anteprime mancanti.} A volte, nel Content Browser, alcune
anteprime non sono renderizzate (la roccia ``non c'è''): Unreal ha avuto
problemi a scattare l'istantanea. Trascinando l'oggetto nel livello, l'engine lo
compila e poi mostra l'anteprima correttamente.

\paragraph{Troppe parentesi.} Durante la dimostrazione si nota anche un difetto
grafico (lo \emph{z-fighting}: due superfici sovrapposte che ``combattono'' per
essere visualizzate), ma il docente si trattiene: ``purtroppo sono una persona
che fa un sacco di parentesi''. Spiegare una cosa gli fa pensare ad altre tre
da spiegare, e cerca di tenere limitato il numero enorme di parentesi che
potrebbe aprire.

% ---------------------------------------------------------------------
\section{Un livello da zero}
\label{gd11:sec:da-zero}

\subsection{Il livello vuoto}

Ora si crea un livello ``proprio partendo da zero, senza usare un template''.
\begin{enumerate}
  \item Nel Content Browser, tornare su \textbf{Contenuto} e lavorare
        \textbf{fuori} dalla cartella dello Starter Content: ``cercate di
        tenere le cose ordinate''.
  \item \textbf{Tasto destro} $\rightarrow$ \textbf{Livello} (\emph{Level}).
        Unreal chiede il nome (il docente lo chiama \emph{NewLevel}).
  \item Aprirlo con un doppio clic.
\end{enumerate}

\paragraph{Il salvataggio.} Aprendo il nuovo livello, Unreal chiede di salvare
ciò che non è stato salvato finora: il livello con cui si è aperto l'engine
risulta \emph{Senza titolo}, cioè non era salvato da nessuna parte. Prima di
cambiare livello, l'editor chiede conferma: ``sicuro che non ti sei dimenticato
qualcosa?''. Va ricordato che \textbf{un nuovo livello non è salvato finché non
lo si salva manualmente} (sezione~\ref{gd8:sec:livelli}).

\paragraph{Il nulla.} ``E abbiamo creato il nostro bellissimo nuovo livello.
Siamo tutti d'accordo che questo è un livello? Siamo tutti appassionati di
domande retoriche?'' L'Outliner è vuoto: a parte le cose tecniche che Unreal
genera per gestirsi, \textbf{nel livello non c'è assolutamente nulla}. Premendo
Play, il vuoto.

\subsection{Un piano, una luce, un cubo}

Bisogna cominciare a metterci qualcosa:
\begin{enumerate}
  \item un \textbf{piano} (dalle forme di base), messo in posizione
        $(0,0,0)$, e poi F per inquadrarlo; lo si allarga, perché il piano di
        partenza è sempre di un metro quadro;
  \item ``che cosa manca? Deselezionando il piano, qual è il problema?''
        \textbf{Manca la luce}. Uno studente propone la luce del cielo
        (\emph{Sky Light}, ``luce atmosferica'' in italiano): inserita da sola,
        non succede niente. Si parte invece dalla \textbf{luce direzionale}:
        quando la si inserisce, Unreal la orienta automaticamente (a circa 46
        gradi) ``e ce l'abbiamo fatta'';
  \item un \textbf{cubo}, ``giusto per avere un punto di riferimento chiaro''.
\end{enumerate}
``Tutto qua?'' No: manca ancora qualcosa.

\subsection{Una parentesi su Lumen}

``Qui devo aprire una parentesi, legata a una roba che ha implementato Unreal.''
Sull'immagine si vede uno \textbf{sfarfallio}, un rumore. Quello è
\textbf{Lumen}.\footnote{Nella trascrizione ``Ruben''.} Chi ha giocato a giochi
come \emph{Clair Obscur: Expedition 33} lo avrà notato in certi angoli: la luce
fa quell'effetto. Lumen è il sistema di \textbf{illuminazione globale dinamica}
di Unreal, ``una imitazione del \emph{ray tracing}'': una luce aggiornata in
tempo reale che considera i \textbf{riflessi delle luci} sulle superfici
circostanti. ``È una finta luce, perché non è una vera luce: è una finzione dei
riflessi della luce su una superficie.'' Infatti, quando la superficie esce
dallo schermo, quegli effetti scompaiono, perché l'engine non sta più calcolando
quei finti riflessi.

\subsection{Perché il lato in ombra è nero}

Guardando il \textbf{lato opposto} del cubo, è \textbf{completamente nero}.
``Se io guardo sotto la mia mano, non è completamente nero. Che cosa sta
succedendo, secondo voi, nel vuoto cosmico?'' Uno studente: la luce si perde nel
vuoto, non si riflette su nessuna superficie.

\paragraph{La luce ambientale.} Esatto: in una situazione di vuoto totale, se
non c'è una superficie che riflette la luce diretta del sole, il lato opposto a
quello illuminato è totalmente nero (``è il motivo per cui la Terra ha la
notte''). Nel mondo reale, invece, ci sono altre fonti di luce oltre al sole. La
prima è l'\textbf{illuminazione ambientale}: la luce che \textbf{si riflette su
qualunque cosa ci circonda}. È il motivo per cui, se si spegnessero le luci e
ogni singolo computer dell'aula, ci sarebbe comunque un po' di luce: è la luce
del sole che rimbalza su tutto ciò che sta fuori ed entra nella stanza.

\paragraph{Il cielo è una fonte di luce.} Seconda cosa: il sole non è l'unica
sorgente; \textbf{anche il cielo è una fonte di luce}. La luce del sole che si
riflette sulla Terra e viene rispedita verso l'alto viene riflessa di nuovo verso
il basso dall'\textbf{atmosfera}. È il motivo per cui il cielo è blu invece di
essere nero (``se vi capiterà mai di andare nello spazio, non lo so cosa
succederà nel 2050--2070'': lo spazio è nero, il cielo della Terra è blu).
L'atmosfera riflette la luce, ed essendo una luce che ci circonda totalmente,
\textbf{contribuisce all'illuminazione ambientale}.

\paragraph{Il cielo, in Unreal.} Ecco perché, trascinando solo una directional
light, Unreal ``dice: bene, vi ho fatto una luce direzionale; e voi guardate e
dite: sì, ok, ma che cos'è questa roba?''. Manca tutto il resto. Bisogna
aggiungere anche la \textbf{Sky Light} (luce atmosferica) e una
\textbf{simulazione del cielo}: la \textbf{Sky Atmosphere} (\emph{atmosfera
cielo}). Inserita quella, il cielo c'è, e c'è il sole.

\subsection{La cartella Lighting: il modo pratico}

Tutto questo si può fare \textbf{a mano}, un elemento alla volta, e nelle slide
c'è il procedimento passo per passo. ``Il vantaggio è che vi posso dare i
consigli pratici da persona che con queste cose ci lavora'': dopo averlo
insegnato un po' di volte, il docente si è reso conto che \textbf{tutti i level
designer}, invece di farlo a mano, \textbf{copiano la cartella}.
\begin{enumerate}
  \item Aprire un livello che ha già la cartella \textbf{Lighting} (per esempio
        quello creato da un template).
  \item Copiarne i \textbf{contenuti}. Attenzione: in Unreal, copiando la
        \emph{cartella} dell'Outliner si copia solo la cartella, non ciò che
        contiene (``questa è la differenza tra lavorare con engine diversi'').
        Bisogna selezionare tutti gli elementi al suo interno.
  \item Tornare al nuovo livello: per cambiare livello velocemente,
        \textbf{File} $\rightarrow$ \textbf{Livelli recenti}.
  \item Assicurarsi che non ci siano già luci direzionali o atmosferiche
        (per non duplicarle) e \textbf{incollare}.
\end{enumerate}
``La cartella è lì apposta per questo.''

\paragraph{Che cosa c'è nella cartella.} Ma bisogna \textbf{sapere che cosa c'è
dentro}, perché regolando quegli elementi si cambia la luce atmosferica del
gioco (tabella~\ref{gd11:tab:lighting}).

\begin{table}[htbp]
\centering
\small
\begin{tabular}{@{}p{4.0cm}p{8.6cm}@{}}
\toprule
\textbf{Elemento} & \textbf{A che cosa serve} \\
\midrule
Directional Light & il sole (sezione~\ref{gd11:sec:luci}) \\
Exponential Height Fog & simula il \textbf{vapore acqueo} nell'atmosfera, che fa
  diventare le cose sempre più sbiadite con la distanza; ``esponenziale''
  perché dipende dalla distanza dal terreno (l'altezza) \\
Volumetric Cloud & le \textbf{nuvole} volumetriche \\
Sky Atmosphere & la simulazione dell'atmosfera: il colore del cielo, il sole,
  il tramonto \\
Sky Light & l'illuminazione ambientale derivata dal cielo \\
Sky Sphere & lo \textbf{sfondo}: di base è bianco (deriva tutto dagli altri
  calcoli), ma ci si può mettere un materiale con un'immagine, come altri
  pianeti o le stelle; è ciò che nei videogiochi si chiama anche
  \textbf{skybox} \\
\bottomrule
\end{tabular}
\caption{Il contenuto della cartella \emph{Lighting} di un livello creato da
template, come descritto a lezione.}
\label{gd11:tab:lighting}
\end{table}

\paragraph{Lo skybox.} Sulla \textbf{sky sphere} il docente si sofferma: è il
vostro sfondo. Nascondendola (con l'icona dell'occhio nell'Outliner) non si nota
niente, perché è totalmente bianca: lo sfondo deriva solo dagli altri calcoli. Ma
se si vogliono altri pianeti sullo sfondo, o le stelle, le si applica un
materiale con un'immagine. Tutto ciò che è \textbf{così lontano dal mondo di
gioco} da poter essere tranquillamente una cosa in 2D, ``tanto non si nota la
differenza''.

\paragraph{Il pavimento.} Fatto tutto questo lavoro (``tutto obbligatorio''), si
prende un \textbf{piano} e lo si allarga in modo che faccia da base del mondo di
gioco. ``Mi state seguendo? Vedo gente che guarda lo schermo senza far nulla.''

% ---------------------------------------------------------------------
\section{Materiali e texture}
\label{gd11:sec:materiali}

\subsection{Assegnare un materiale}

Selezionando il piano, in basso a destra nei \emph{Details} c'è una casella
\textbf{Materiali}: da lì si sceglie il materiale da applicare. Il docente vuole
un pavimento di pietre: sceglie \textbf{Cobblestone} (ciottolato), poi prova
altro.
\begin{itemize}
  \item Si può \textbf{scorrere} l'elenco, oppure \textbf{cercare} per nome:
        ``voglio che la mia base sia erba'': si scrive \emph{grass}, e si trova
        l'erba.
  \item Oppure si apre lo Starter Content, cartella \textbf{Materials}, si
        prende il materiale e lo si \textbf{trascina} sull'oggetto: lo applica
        automaticamente (per esempio i mattoni).
  \item Selezionando più oggetti insieme, si può cambiare il materiale a
        tutti in contemporanea.
\end{itemize}

\subsection{Il tiling e lo stretch}

``Tenete presente che questo è tarato per essere un metro per un metro.'' È il
discorso visto con le spline (sezione~\ref{gd11:sec:spline}): allargando o
riducendo un oggetto, \textbf{le dimensioni della texture rispetto all'oggetto
non cambiano}, perché le dimensioni locali sono sempre le stesse. Quindi, se lo
si allunga, \textbf{la texture si stira}.

\paragraph{Il pavimento a piastrelle.} Per fare un pavimento senza che sia
``tutto stretchato e brutto da vedere'', per ora:
\begin{enumerate}
  \item scegliere una \textbf{scala di base} per il piano: 1, oppure 2 se va
        bene che sia leggermente stirata;
  \item impostare la \textbf{griglia a 100} (un metro per un metro,
        sezione~\ref{gd11:sec:snapping});
  \item fare un po' di \textbf{copia e incolla}, e affiancare i pezzi.
\end{enumerate}
È un metodo da principianti, ma è voluto: ``per il momento in cui siamo noi,
lavoriamo in questo modo, invece di guardare quindici strumenti diversi in una
volta sola''.

\paragraph{Le giunture.} Queste texture sono pensate per essere
\textbf{ripetibili}, ma non tutte alla perfezione: tra un pezzo e l'altro si
possono vedere i segni. Se dà fastidio, si cerca un materiale con un
\textbf{tiling} migliore. Ma dalla distanza ``non si nota tanto'', e in questa
fase va bene: ``siamo agli inizi, non preoccupiamoci di UV e di tutte le altre
cose''.\footnote{Durante la dimostrazione due piani che dovrebbero avere lo
stesso materiale risultano avere materiali diversi (``che roba strana''), e il
computer si blocca per qualche secondo; il docente non approfondisce.}

\subsection{Texture e materiale non sono la stessa cosa}

``Una nota importante: \textbf{materiali e texture sono due cose diverse}.''
\begin{description}
  \item[Texture] è \textbf{un componente} di un materiale: il \textbf{file 2D},
        l'immagine (per esempio un PNG).
  \item[Materiale] è \textbf{un insieme di texture più un sacco di altri
        calcoli}, inclusi quelli che fanno sì che la texture \textbf{finga di
        essere in rilievo}.
\end{description}
Sui mattoni si vede che i mattoni ``sporgono''; in realtà non sporgono affatto,
la superficie è completamente piatta: è Unreal che, calcolandoli, finge che
sporgano nel calcolare l'immagine da mostrare a schermo.\footnote{Il docente
accenna, più avanti, alle \emph{normali} (``quella roba di profondità che vi ho
fatto vedere prima, che sono le normali''): l'effetto di rilievo si ottiene con
una \emph{normal map}, una texture che indica all'engine come la luce deve
comportarsi su ogni punto della superficie. Se ne parlerà in un'altra
lezione.} I materiali si trovano nella casella vista prima; e nelle slide c'è la
spiegazione, un passaggio alla volta, di tutti gli elementi visti
(\emph{Exponential Height Fog} compresa).

% ---------------------------------------------------------------------
\section{Le modalità di visualizzazione}
\label{gd11:sec:viewmode}

``Quanto tempo abbiamo? Ok, ce l'abbiamo.'' In alto nella viewport c'è una
casella con scritto \textbf{Illuminato}, in inglese \textbf{Lit} (``io odio come
è tradotto in italiano: è molto più comodo fare ricerche se l'avete in inglese,
ma scegliete voi''). È la \textbf{modalità di visualizzazione}, o \emph{View
Mode}. Si sceglie a seconda di ciò che serve vedere (tabella~\ref{gd11:tab:view}).

\begin{table}[htbp]
\centering
\small
\begin{tabular}{@{}p{3.6cm}p{9.0cm}@{}}
\toprule
\textbf{Modalità} & \textbf{Che cosa mostra, e quando usarla} \\
\midrule
Illuminato (\emph{Lit}) & \textbf{tutti i calcoli} che danno l'immagine finale:
  ombre, luci, post-processing. È la modalità ``totale''. \\
\addlinespace
Wireframe & i \textbf{poligoni} di cui è composto ogni oggetto: si vedono i
  triangoli (alcuni poligoni ne raggruppano più d'uno). Utile per capire com'è
  fatta una mesh. \\
\addlinespace
Non illuminato (\emph{Unlit}) & praticamente \textbf{solo le texture}, ``come
  facevano i vecchi videogiochi, tipo \emph{Quake}, \emph{Doom}'': nessun calcolo
  di illuminazione, nemmeno l'effetto di rilievo delle normali. Utile per
  lavorare quando si piazzano pareti e oggetti, o se Unreal è pesante sul
  proprio computer. \\
\bottomrule
\end{tabular}
\caption{Le tre modalità di visualizzazione su cui fermarsi per ora.}
\label{gd11:tab:view}
\end{table}

\paragraph{Wireframe.} Su una scena con pochi oggetti, il wireframe non mostra
molto; trascinando qualcosa di più complesso dallo Starter Content (una sedia
tra i \emph{Props}) si vedono i \textbf{triangoli}, come la testa della slide
sulle mesh (sezione~\ref{gd11:sec:mesh}).

\paragraph{Lavorare in non illuminato.} Se si stanno solo mettendo pareti e
oggetti in una stanza, \textbf{non serve che il computer calcoli sempre tutte le
ombre}: si può lavorare tranquillamente in \emph{Unlit} e passare a \emph{Lit}
solo quando si lavora sulle luci o si vuole vedere il risultato. Se Unreal fa
fatica sul proprio computer, è un modo per semplificarsi la vita. Ci sono poi
anche le \textbf{impostazioni di scalabilità}, che si vedranno in un altro
momento.

\paragraph{Molte altre.} Ci sono tantissime altre modalità di visualizzazione,
che servono per esempio a vedere se c'è qualche problema; per ora si resta su
illuminato, non illuminato e wireframe, ``se volete dare un'occhiata e vedere
cosa succede dietro le quinte''.

% ---------------------------------------------------------------------
\section{Il post process volume}
\label{gd11:sec:postprocess}

``Un'ultima cosa che voglio dire, e questa è per voi, se volete divertirvi'',
``se volete proprio esagerare''. Tra gli elementi da aggiungere al livello c'è il
\textbf{Post Process Volume} (in italiano \emph{volume di post processing}): un
cubo che permette di aggiungere tutte le opzioni di \textbf{post-processing}.

\paragraph{Che cos'è il post-processing.} Sono tutte le impostazioni dei
\textbf{calcoli che vengono fatti sull'immagine una volta che è stata
renderizzata}. ``Detto così è molto tecnico.'' In pratica:
\begin{itemize}
  \item il \textbf{motion blur};
  \item il \textbf{bloom}: le cose particolarmente illuminate espandono la
        loro luce oltre i propri confini;
  \item la \textbf{profondità di campo} (\emph{depth of field}): se una cosa in
        primo piano è a fuoco, lo sfondo non lo è;
  \item la \textbf{grana della pellicola} (\emph{film grain}), la
        \textbf{correzione del colore}, il \textbf{lens flare},
        l'\textbf{aberrazione cromatica}, e così via.
\end{itemize}
Tutto è calcolato dopo che la scena è stata renderizzata. Sono effetti
abbastanza complessi da usare, perché ci sono tanti valori da valutare; non è
``che clicchi e lo fai'', richiede un po' di studio. Ma se si vuole ``l'effetto
videogioco super figo'', il post process volume è lo strumento. Ci sarà una
lezione dedicata; per ora il docente lo lascia esplorare da soli, ``perché c'è
sempre qualcuno che vuole strafare e si prende tantissimo bene''.

\paragraph{Che cos'è un volume.} Il post process volume è un
\textbf{volume}, un altro elemento di Unreal: un oggetto nel gioco che
\textbf{fa succedere qualcosa solo se altre cose sono al suo interno}. Il cubo
delimita un'area: le impostazioni di post-processing si applicano solo se
\textbf{la telecamera è dentro il volume}. Serve, per esempio, a cambiare il
post-processing da una zona all'altra del livello.

\paragraph{Infinite Extent.} L'unica cosa da ricordare (anche questa segnata
sulla slide) è la casella \textbf{Infinite Extent (Unbound)}: se la si attiva,
il volume si applica \textbf{a tutto il livello}, ovunque sia la telecamera. Il
motivo per cui all'inizio della dimostrazione non si vedeva niente era proprio
che non era attiva. Una volta attivata, alzando l'intensità del \textbf{bloom},
la luce diventa sempre più forte, soprattutto quella riflessa; scrivendo a mano
un valore oltre il limite (50), ``tutto si spara, flasha'', perché il bloom è al
massimo.

\paragraph{Una svista.} Il docente prova anche la profondità di campo, che è più
complessa da impostare perché vanno regolati tutti i parametri, e apre per
sbaglio la sezione per il \emph{mobile}: ``scusate, ho detto una cavolata. La
mobile non funziona in questo, perché non stiamo facendo un gioco per
cellulare''. Rimandato alla lezione dedicata.

% ---------------------------------------------------------------------
\section{Il task: una stanza, qualunque stanza}
\label{gd11:sec:task}

``Abbiamo coperto un sacco di roba oggi. E queste sono tutte lezioni di pratica
iniziale.'' Quindi il task per la lezione successiva è, nella forma, lo stesso
della volta precedente, ma con una differenza importante: non si ricopia più la
propria stanza.

\begin{quote}
\textbf{Creare una stanza}, ``qualcosa, qualunque cosa, che sia una stanza''.
Possono essere anche più di una, ma \textbf{minimo una}. Usare \textbf{gli asset
e i primitivi} a disposizione: le forme di base usate finora, i
\textbf{materiali} visti oggi e le varie cose disponibili nello
\textbf{Starter Content} (prop e così via). Consegna: \textbf{almeno tre
screenshot}, nella cartella dei task, per la settimana successiva.
\end{quote}

\paragraph{Un task creativo.} ``Questa è una task creativa, quindi divertitevi:
potete scegliere voi che cosa fate.'' L'unico vincolo è che siano
\textbf{stanze}: quest'anno si è deciso di concentrarsi sugli \textbf{interni},
per evitare troppa dispersione. Dentro le stanze, però, ``potete farci quel che
volete, non vi do nessuna limitazione''. Lo Starter Content non ha molti prop: è
più utile per i \textbf{materiali} (erba, legno, pietra, mattoni, ``tutte quelle
robe che possono coprire le vostre cose di base'').

\paragraph{Tre screenshot: un minimo.} Uno studente chiede se le tre foto siano
un obbligo o un minimo. ``Per me è un minimo'': con una sola foto si vede solo
una parete. Si può mandarne di più (``diciamo minimo tre''). Altri docenti sono
più rigidi, ma lui non vuole limitare chi ha idee particolari, chi vuole
mostrare angoli specifici della stanza o preferisce mostrare di più. Un'altra
domanda riguarda se ogni oggetto debba avere necessariamente un materiale: si
può anche lasciare un oggetto senza, se ha senso.

\paragraph{Guardate gli altri videogiochi.} ``Una nota importante. Io sono uno
che ci tiene a dire: guardate gli altri videogiochi come \textbf{punto di
riferimento}, cercate di capire che cosa fanno gli altri.'' Nel mondo dei
videogiochi professionali non si guarda solo il proprio lavoro: ci si guarda
attorno.

\paragraph{Ma siete all'inizio.} Allo stesso tempo, ``è chiaro che, se vi do
questo task, la settimana prossima non mi aspetto che mi portiate una camera da
letto al livello di \emph{GTA 6}'', anche perché si hanno a disposizione solo
lo Starter Content e le forme di base. Non si può fare roba iperrealistica.
``Metteteci un pochino di impegno, fate un po' quello che volete, ma non
impazzite a fare cambiamenti super minuscoli'': si è studenti appena
all'inizio. Il principio:
\begin{quote}
Usate gli strumenti che avete \textbf{nello scopo degli strumenti che avete}.
Poi, piano piano, man mano che avremo sempre più strumenti a disposizione,
porteremo sempre più in alto il livello del nostro lavoro.
\end{quote}
Tutti questi primi esercizi servono a \textbf{prendere la mano con Unreal}: ``mi
interessa molto di più che vedere una roba artisticamente perfetta; voglio vedere
che ci mettete l'impegno''.

\paragraph{Approfondimento: dalla stanza copiata alla stanza inventata.} Il
passaggio dal primo al secondo task non è casuale. Ricopiare la propria stanza
(sezione~\ref{gd8:sec:task}) allena la mano e il senso della scala; inventare una
stanza obbliga a fare le scelte che finora erano già date: che funzione ha lo
spazio, chi ci vive, come lo si \emph{legge}. È lì che entrano in gioco le
osservazioni della revisione: il \textbf{pizzico di caos} che rende uno spazio
vissuto (le scatole sghembe di Federica), il rischio dello \textbf{showroom}
(la stanza di Alberto), il criterio del \textbf{verosimile} (il letto di
Francesco). Totten insiste su un'idea simile quando parla di \emph{storytelling
con asset modulari}: anche pochi pezzi standard, scelti e disposti con
intenzione, possono raccontare chi abita uno spazio.

% ---------------------------------------------------------------------
\section{Le domande e la chiusura}
\label{gd11:sec:chiusura}

\paragraph{Le ultime domande.} Come sempre, gli ultimi minuti sono per domande e
curiosità, e per chi vuole cominciare il task in classe. Uno studente non riesce a far comparire lo Starter Content: il docente gli
rispiega che la
cartella va trascinata \textbf{dentro la cartella del progetto aperta in
Windows}, non nel Content Browser, perché ``non state importando nuovi file: mi
sono stati passati dei file che sono già stati importati''. Poi si chiude e si
riapre Unreal, che li trova da solo (sezione~\ref{gd11:sec:starter}). ``Lo so,
era molto più comodo quando lo potevi selezionare con un clic.'' Il docente
promette di studiare meglio la questione dopo la lezione: ``veramente antipatico,
ma lo dovrò studiare''.

\paragraph{Le slide.} ``Come continuo a ripetere, le slide sono dentro la
cartella.''

\paragraph{Le presenze.} Ci sono ancora due o tre minuti, ma ``possiamo anche
chiudere qui''. Mattia, da casa, conferma che ha avuto problemi solo con il
computer. E poi l'ultima cosa, che il docente stava per dimenticare ``con tutto
quello che è successo'': le \textbf{presenze}. Attadia, Etiopia, Ferrara, Maestri
(da casa), Muratore, Ronchi, Rossaro e Vergata ci sono; Maracino è assente.
``Alla settimana prossima. Buon lavoro, arrivederci.''

% ---------------------------------------------------------------------
\section*{Punti chiave del capitolo}
\addcontentsline{toc}{section}{Punti chiave del capitolo}

\begin{itemize}
  \item Il primo anno, ogni lezione di level design comincia con la
        \textbf{revisione dei task} in classe: per vedere il livello generale
        e correggere una volta sola gli errori comuni. Nella cartella dei task
        non si cancella: si scrive ``da cancellare''. Screenshot con
        \textbf{Win+Shift+S}.
  \item Dalla revisione: \textbf{un pizzico di caos} (oggetti sghembi) rende
        gli ambienti vissuti; una stanza vuota sembra uno showroom; i
        videogiochi puntano al \textbf{verosimile}, non al realistico. La scala
        è un moltiplicatore: scala 5 su un cubo da un metro fa cinque metri.
  \item Il level design ha due macro-campi: \textbf{artistico} (luci,
        \emph{enrichment}, post-processing, VFX) e \textbf{di design}
        (\emph{encounter design}, blockout, funzione degli spazi). Negli studi
        grandi ci si specializza, e nei test di assunzione spesso si sceglie
        tra test artistico e test di gameplay. Il primo anno è sul lato
        artistico.
  \item Il \textbf{blockout} (\emph{grey box}, \emph{white box}, \emph{block
        mesh}: sinonimi) è lo sketch a matita del livello: forme semplici per
        arrivare \textbf{il prima possibile a qualcosa di giocabile}, testare
        idee e rispondere a domande. Segnala l'\textbf{intento} (la finestra
        rossa di \emph{Call of Duty}: ``da lì possono spararti'').
  \item Prima il gameplay, la funzione e la forma degli spazi, poi
        l'estetizzazione: \textbf{il gameplay regna sovrano}. Ma la veste
        artistica è decisa prima, e si lavora in parallelo, non a cascata.
  \item Cinque strumenti, non esclusivi: \textbf{primitivi} (semplici ma
        limitati), \textbf{kit modulari} (velocissimi ma rigidi; dungeon di
        \emph{Skyrim} ed \emph{Elden Ring}, roguelike), \textbf{brush e
        modellazione} (qualunque forma, curva di apprendimento; Hammer di
        Valve, Blender per Unity), \textbf{scultura del terreno} (colline),
        \textbf{spline} (strade e cavi: copie di una mesh lungo una curva,
        per evitare lo stretch). Si sceglie lo strumento più veloce per lo
        scopo.
  \item \textbf{Asset} = file nel progetto; \textbf{actor} = istanza nel
        livello. Gli asset dell'engine non si modificano.
  \item \textbf{Coordinate globali}: dall'origine $(0,0,0)$ del mondo (non è lo
        spawn); in Unreal $z$ è l'alto. \textbf{Coordinate locali}: dall'origine
        dell'oggetto o del genitore. Nell'editor: spazio del mondo, locale,
        genitore; la scala è sempre locale.
  \item Un \textbf{figlio} eredita tutte le trasformazioni del genitore
        (anche la scala, distanze comprese); le gerarchie si propagano verso il
        basso. Un \textbf{gruppo} (Ctrl+G) seleziona e muove insieme oggetti
        indipendenti. L'asterisco segnala modifiche non salvate.
  \item \textbf{Ctrl+Z annulla anche le selezioni}: cliccare è un'azione.
  \item \textbf{Snapping}: griglia in centimetri (10 di default; 100 per il
        cubo da un metro), rotazione a 10 gradi; spegnerlo per oggetti
        naturali. Valori a mano e operazioni nei campi. Ruotare
        \emph{sull'asse} $z$ significa ruotare \emph{attorno} alla verticale.
  \item Content Browser: Ctrl+Spazio. Lingua: Impostazioni dell'editor. Il
        \textbf{pivot} è l'origine delle coordinate di un actor: Alt + tasto
        centrale lo sposta temporaneamente (anche fuori dall'oggetto). G
        nasconde i gizmo, F inquadra l'oggetto selezionato.
  \item Una \textbf{mesh} è un insieme di punti e segmenti che formano
        triangoli; \textbf{mesh e collisioni sono cose diverse}: la collisione
        non è automatica, va controllata.
  \item Luci: \textbf{point} (sfera, raggio di attenuazione, source radius),
        \textbf{spot} (cono), \textbf{rect} (superficie),
        \textbf{directional} (il sole: conta la rotazione, una sola per
        livello). Il modello della lampadina non emette luce. Alcuni campi
        accettano valori oltre il limite del cursore.
  \item Lo \textbf{Starter Content} nella 5.8 va aggiunto a mano: decomprimere,
        trascinare la cartella nella cartella \texttt{Content} del progetto,
        riaprire Unreal.
  \item Un livello da zero ha bisogno di sole, \textbf{luce ambientale} e
        \textbf{cielo} (Sky Light, Sky Atmosphere, fog, nuvole, sky sphere): in
        pratica si copiano i contenuti della cartella \emph{Lighting}. Lumen
        simula i riflessi della luce in tempo reale.
  \item \textbf{Texture} = immagine 2D; \textbf{materiale} = texture + calcoli
        (anche il finto rilievo). Scalare un oggetto stira la texture: per ora,
        piastrelle con griglia a 100.
  \item View mode: \textbf{Lit}, \textbf{Unlit} (più leggero), \textbf{Wireframe}.
        Il \textbf{Post Process Volume} agisce solo dentro il volume, a meno
        di \textbf{Infinite Extent}.
  \item Task: \textbf{una stanza inventata} (almeno una), con primitivi,
        Starter Content e materiali; almeno tre screenshot. Guardate gli altri
        giochi, ma usate gli strumenti nello scopo degli strumenti che avete.
\end{itemize}
